\documentclass[../../main.tex]{subfiles}
\graphicspath{{\subfix{../../image/}}}

\begin{document}

\section{Sobolev不等式}\label{section:ch05-Sobolev不等式}

本节的目标是发现各种Sobolev空间之间的嵌入。此处关键的分析工具将是某些所谓的“Sobolev型不等式”，我们下面先对光滑函数证明这些不等式。由于——正如我们在\cref{section:ch05-逼近} 中所见——光滑函数稠密，这些不等式随后将为各个相关Sobolev空间中的任意函数建立估计。

为清晰起见，我们先只考虑Sobolev空间 $W^{1,p}(U)$，并提出如下基本问题：若函数 $u$ 属于 $W^{1,p}(U)$，它是否自动属于某些其他空间？答案是“是”，但究竟属于哪些其他空间取决于
\begin{gather}
1\leqslant p<n, \label{align:ch05-H-1表示系数}\\
p=n, \label{align:ch05-情形2}\\
n<p\leqslant\infty. \label{align:ch05-情形3}
\end{gather}
我们在\cref{subsection:ch05-GNS} 研究情形~\eqref{align:ch05-H-1表示系数}，在\cref{subsection:ch05-Morrey} 研究情形~\eqref{align:ch05-情形3}，而边界情形~\eqref{align:ch05-情形2} 留待\cref{subsection:ch05-边界情形} 讨论。

\subsection{Gagliardo--Nirenberg--Sobolev 不等式}\label{subsection:ch05-GNS}

本节设
\begin{align}
1\leqslant p<n
\end{align}
并先追问能否建立形如
\begin{align}
\|u\|_{L^{q}(\mathbb{R}^{n})}\leqslant C\|Du\|_{L^{p}(\mathbb{R}^{n})}
\label{align:ch05-H1估计}
\end{align}
的估计，其中常数 $C>0$、$1\leqslant q<\infty$，对一切 $u\in C^{\infty}_{c}(\mathbb{R}^{n})$ 成立。要点在于常数 $C$ 与 $q$ 都不应依赖于 $u$。

\textbf{动机。} 我们首先证明：若任何形如~\eqref{align:ch05-H1估计} 的不等式成立，则数 $q$ 不能任意，而事实上必须具有非常特定的形式。为此取函数 $u\in C^{\infty}_{c}(\mathbb{R}^{n})$，$u\not\equiv 0$，并对 $\lambda>0$ 定义伸缩函数
\begin{align*}
u_{\lambda}(x)\triangleq u(\lambda x)\qquad (x\in\mathbb{R}^{n}).
\end{align*}
把式~\eqref{align:ch05-H1估计} 应用于 $u_{\lambda}$，得
\begin{align}
\|u_{\lambda}\|_{L^{q}(\mathbb{R}^{n})}\leqslant C\|Du_{\lambda}\|_{L^{p}(\mathbb{R}^{n})}.
\end{align}
现
\begin{align*}
\int_{\mathbb{R}^{n}}|u_{\lambda}|^{q}\,\mathrm{d}x
=\int_{\mathbb{R}^{n}}|u(\lambda x)|^{q}\,\mathrm{d}x
=\frac{1}{\lambda^{n}}\int_{\mathbb{R}^{n}}|u(y)|^{q}\,\mathrm{d}y,
\end{align*}
且
\begin{align*}
\int_{\mathbb{R}^{n}}|Du_{\lambda}|^{p}\,\mathrm{d}x
=\lambda^{p}\int_{\mathbb{R}^{n}}|Du(\lambda x)|^{p}\,\mathrm{d}x
=\frac{\lambda^{p}}{\lambda^{n}}\int_{\mathbb{R}^{n}}|Du(y)|^{p}\,\mathrm{d}y.
\end{align*}
把这些等式代入式~\eqref{align:ch05-H1估计}，得
\begin{align*}
\frac{1}{\lambda^{\frac{n}{q}}}\|u\|_{L^{q}(\mathbb{R}^{n})}
\leqslant C\frac{\lambda}{\lambda^{\frac{n}{p}}}\|Du\|_{L^{p}(\mathbb{R}^{n})},
\end{align*}
故
\begin{align}
\|u\|_{L^{q}(\mathbb{R}^{n})}\leqslant C\lambda^{1-\frac{n}{p}+\frac{n}{q}}\|Du\|_{L^{p}(\mathbb{R}^{n})}.
\end{align}
但若 $1-\dfrac{n}{p}+\dfrac{n}{q}\neq 0$，则在该式中令 $\lambda$ 趋于 $0$ 或 $\infty$ 便得矛盾。因此若所期望的不等式~\eqref{align:ch05-H1估计} 确实成立，我们必须有 $1-\dfrac{n}{p}+\dfrac{n}{q}=0$；于是
\begin{align*}
\frac{1}{q}=\frac{1}{p}-\frac{1}{n},\qquad q=\frac{np}{n-p}.
\end{align*}
这一观察促使我们给出

\begin{definition}\label{definition:ch05-Sobolev共轭}
若 $1\leqslant p<n$，$p$ 的\textbf{Sobolev共轭}为
\begin{align}
p^{*}\triangleq\frac{np}{n-p}.
\end{align}
注意
\begin{align}
\frac{1}{p^{*}}=\frac{1}{p}-\frac{1}{n},\qquad p^{*}>p.
\end{align}
\end{definition}
\begin{remark}
上述伸缩分析表明估计~\eqref{align:ch05-H1估计} 只可能对 $q=p^{*}$ 成立。
\end{remark}

\begin{theorem}[Gagliardo--Nirenberg--Sobolev不等式]\label{theorem:Gagliardo--Nirenberg--Sobolev不等式}
设 $1\leqslant p<n$。则存在一个只依赖于 $p$ 与 $n$的常数 $C=\frac{p(n-1)}{n-p}$，使对一切 $u\in C^{1}_{c}(\mathbb{R}^{n})$ 有
\begin{align}
\|u\|_{L^{p^{*}}(\mathbb{R}^{n})}\leqslant C\|Du\|_{L^{p}(\mathbb{R}^{n})}.
\label{align:ch05-GNS}
\end{align}
\end{theorem}
\begin{remark}
现式~\eqref{align:ch05-GNS} 确实需要 $u$ 具有紧支集才成立，正如 $u\equiv 1$ 的例子所示。但值得注意的是，此处的常数完全不依赖于 $u$ 的支集大小.
\end{remark}
\begin{proof}

1. 先设 $p=1$。由于 $u$ 有紧支集，故当$|x|$充分大时,有$u(x)=0$.对每个 $i=1,\cdots,n$ 及 $x\in\mathbb{R}^{n}$，由微积分基本定理有
\begin{align*}
u(x)=u(x)-u(-\infty)=\int_{-\infty}^{x_{i}}u_{x_{i}}(x_{1},\cdots,x_{i-1},y_{i},x_{i+1},\cdots,x_{n})\,\mathrm{d}y_{i};
\end{align*}
因而
\begin{align*}
|u(x)|\leqslant\int_{-\infty}^{\infty}|Du(x_{1},\cdots,y_{i},\cdots,x_{n})|\,\mathrm{d}y_{i}\quad (i=1,\cdots,n).
\end{align*}
于是
\begin{align*}
|u(x)|^{\frac{n}{n-1}}
\leqslant\prod_{i=1}^{n}\left(\int_{-\infty}^{\infty}|Du(x_{1},\cdots,y_{i},\cdots,x_{n})|\,\mathrm{d}y_{i}\right)^{\frac{1}{n-1}}.
\end{align*}
将此不等式对 $x_{1}$ 积分得
\begin{align*}
\int_{-\infty}^{\infty}|u|^{\frac{n}{n-1}}\,\mathrm{d}x_{1}
&\leqslant\int_{-\infty}^{\infty}\prod_{i=1}^{n}\left(\int_{-\infty}^{\infty}|Du|\,\mathrm{d}y_{i}\right)^{\frac{1}{n-1}}\mathrm{d}x_{1}\\
&=\left(\int_{-\infty}^{\infty}|Du|\,\mathrm{d}y_{1}\right)^{\frac{1}{n-1}}
\int_{-\infty}^{\infty}\prod_{i=2}^{n}\left(\int_{-\infty}^{\infty}|Du|\,\mathrm{d}y_{i}\right)^{\frac{1}{n-1}}\mathrm{d}x_{1}\\
&\leqslant\left(\int_{-\infty}^{\infty}|Du|\,\mathrm{d}y_{1}\right)^{\frac{1}{n-1}}
\prod_{i=2}^{n}\left(\int_{-\infty}^{\infty}\int_{-\infty}^{\infty}|Du|\,\mathrm{d}x_{1}\,\mathrm{d}y_{i}\right)^{\frac{1}{n-1}},
\end{align*}
最后一个不等式来自\hyperref[theorem:广义Holder不等式]{广义Hölder不等式}。再将其对 $x_{2}$ 积分，并再次应用\hyperref[theorem:广义Holder不等式]{广义Hölder不等式}，如此继续对 $x_{3},\cdots,x_{n}$ 积分，最终得到
\begin{align}
\int_{\mathbb{R}^{n}}|u|^{\frac{n}{n-1}}\,\mathrm{d}x
\leqslant\prod_{i=1}^{n}\left(\int_{-\infty}^{\infty}\cdots\int_{-\infty}^{\infty}|Du|\,\mathrm{d}x_{1}\cdots\mathrm{d}y_{i}\cdots\mathrm{d}x_{n}\right)^{\frac{1}{n-1}}
=\left(\int_{\mathbb{R}^{n}}|Du|\,\mathrm{d}x\right)^{\frac{n}{n-1}}.\label{eq:JIFOJEIOJIOF32423rIDFJ}
\end{align}
这就是 $p=1$ 时的估计式\eqref{align:ch05-GNS}。

2. 现考虑 $1<p<n$ 的情形。在\eqref{eq:JIFOJEIOJIOF32423rIDFJ}式中将函数$u$改为$v\triangleq|u|^{\gamma}$，其中 $\gamma>1$ 待定。则由\hyperref[Real Analysis-theorem:实变函数--Hölder不等式]{Hölder不等式}可得
\begin{align}\label{eq:jfoiwj3::ijo2jjj}
\left(\int_{\mathbb{R}^{n}}|u|^{\frac{\gamma n}{n-1}}\,\mathrm{d}x\right)^{\frac{n-1}{n}}
\leqslant\int_{\mathbb{R}^{n}}\left|D|u|^{\gamma}\right|\,\mathrm{d}x
=\gamma\int_{\mathbb{R}^{n}}|u|^{\gamma-1}|Du|\,\mathrm{d}x
\leqslant\gamma\left(\int_{\mathbb{R}^{n}}|u|^{(\gamma-1)\frac{p}{p-1}}\,\mathrm{d}x\right)^{\frac{p-1}{p}}
\left(\int_{\mathbb{R}^{n}}|Du|^{p}\,\mathrm{d}x\right)^{\frac{1}{p}}.
\end{align}
我们选 $\gamma$ 使 $\dfrac{\gamma n}{n-1}=(\gamma-1)\dfrac{p}{p-1}$。即令
\begin{align*}
\gamma\triangleq\frac{p(n-1)}{n-p}>1,
\end{align*}
此时
\begin{align*}
\frac{\gamma n}{n-1}=(\gamma-1)\frac{p}{p-1}=\frac{np}{n-p}=p^{*}.
\end{align*}
于是估计式\eqref{eq:jfoiwj3::ijo2jjj}变为
\begin{align*}
\left(\int_{\mathbb{R}^{n}}|u|^{p^{*}}\,\mathrm{d}x\right)^{\frac{1}{p^{*}}}
\leqslant C\left(\int_{\mathbb{R}^{n}}|Du|^{p}\,\mathrm{d}x\right)^{\frac{1}{p}}.
\end{align*}
\end{proof}

\begin{corollary}\label{corollary:Sobolev嵌入Lqq大于n}
设 $n \geqslant 2$。则对一切 $q \in [n, +\infty)$，有
\begin{align*}
W^{1,n}(\mathbb{R}^n) \subseteq L^q(\mathbb{R}^n),
\end{align*}
并且存在常数 $C$ 只依赖于 $n$ 与 $q$，使得
\begin{align*}
\|u\|_{L^q(\mathbb{R}^n)} \leqslant C\|u\|_{W^{1,n}(\mathbb{R}^n)}.
\end{align*}
\end{corollary}
\begin{proof}

先设 $u \in C_c^1(\mathbb{R}^n)$。对 $|u|^m$（其中 $m \geqslant 1$）应用 $p=1$ 时的 \hyperref[theorem:Gagliardo--Nirenberg--Sobolev不等式]{Gagliardo--Nirenberg--Sobolev不等式}，得到
\begin{align*}
\||u|^m\|_{L^{\frac{n}{n-1}}(\mathbb{R}^n)} \leqslant C_1\|D(|u|^m)\|_{L^1(\mathbb{R}^n)}.
\end{align*}
由链式法则得$D(|u|^m) = m|u|^{m-1}\operatorname{sgn}(u)Du,$
因此 $|D(|u|^m)| = m|u|^{m-1}|Du|$。代入上式得
\begin{align}\label{eq::brezis-limit-first1}
\|u\|_{L^{\frac{mn}{n-1}}(\mathbb{R}^n)}^m \leqslant C_1 m \int_{\mathbb{R}^n} |u|^{m-1}|Du| \,\mathrm{d}x.
\end{align}
对右端积分使用\hyperref[Real Analysis-theorem:实变函数--Hölder不等式]{Hölder不等式}，得
\begin{align*}
\int_{\mathbb{R}^n} |u|^{m-1}|Du| \,\mathrm{d}x
\leqslant \left( \int_{\mathbb{R}^n} (|u|^{m-1})^{\frac{n}{n-1}} \,\mathrm{d}x \right)^{\frac{n-1}{n}} \left( \int_{\mathbb{R}^n} |Du|^n \,\mathrm{d}x \right)^{\frac{1}{n}}= \|u\|_{L^{\frac{(m-1)n}{n-1}}(\mathbb{R}^n)}^{m-1} \|Du\|_{L^n(\mathbb{R}^n)}.
\end{align*}
代入\eqref{eq::brezis-limit-first1}式，得到
\begin{align}\label{eq::brezis-limit-first}
\|u\|_{L^{\frac{mn}{n-1}}(\mathbb{R}^n)}^m \leqslant C_1 m\|u\|_{L^{\frac{(m-1)n}{n-1}}(\mathbb{R}^n)}^{m-1}\|Du\|_{L^n(\mathbb{R}^n)}.
\end{align}
令
\begin{align*}
A\triangleq\|u\|_{L^{\frac{mn}{n-1}}(\mathbb{R}^n)},\quad
B\triangleq\|u\|_{L^{\frac{(m-1)n}{n-1}}(\mathbb{R}^n)},\quad
D\triangleq\|Du\|_{L^n(\mathbb{R}^n)}.
\end{align*}
则由 \eqref{eq::brezis-limit-first} 式得 $A^m \leqslant C_1 m B^{m-1}D$。两边开 $m$ 次方得
\begin{align*}
A \leqslant (C_1 m)^{\frac{1}{m}} B^{\frac{m-1}{m}} D^{\frac{1}{m}}.
\end{align*}
由\hyperref[Basis of Analytics-theorem:Young不等式]{Young不等式}，对于指数 $p=m$ 与 $p'=\frac{m}{m-1}$，有
\begin{align*}
B^{\frac{m-1}{m}} D^{\frac{1}{m}} \leqslant \frac{m-1}{m}B + \frac{1}{m}D \leqslant B + D.
\end{align*}
因此
\begin{align*}
A \leqslant (C_1 m)^{\frac{1}{m}} (B + D).
\end{align*}
记 $C_m\triangleq(C_1 m)^{\frac{1}{m}}$，它只依赖于 $m$ 与 $n$，于是
\begin{align}\label{eq::brezis-limit-young}
\|u\|_{L^{\frac{mn}{n-1}}(\mathbb{R}^n)} \leqslant C_m\left(\|u\|_{L^{\frac{(m-1)n}{n-1}}(\mathbb{R}^n)} + \|Du\|_{L^n(\mathbb{R}^n)}\right).
\end{align}

在 \eqref{eq::brezis-limit-young} 式中取 $m=n$，则左边为 $\|u\|_{L^{\frac{n^2}{n-1}}(\mathbb{R}^n)}$，右边括号内为 $\|u\|_{L^n(\mathbb{R}^n)} + \|Du\|_{L^n(\mathbb{R}^n)}$，而后者被 $\|u\|_{W^{1,n}(\mathbb{R}^n)}$ 控制。因此
\begin{align}\label{eq::brezis-limit-mn}
\|u\|_{L^{\frac{n^2}{n-1}}(\mathbb{R}^n)} \leqslant C\|u\|_{W^{1,n}(\mathbb{R}^n)}.
\end{align}
由\hyperref[theorem:广义Holder不等式]{广义Hölder不等式}，对一切 $q$ 满足 $n \leqslant q \leqslant \frac{n^2}{n-1}$，存在 $\theta \in [0,1]$ 使得
\begin{align*}
\frac{1}{q} = \frac{\theta}{n} + \frac{1-\theta}{\frac{n^2}{n-1}} \iff 1= \frac{\theta q}{n} + \frac{(1-\theta)q}{\frac{n^2}{n-1}}.
\end{align*}
于是结合\eqref{eq::brezis-limit-mn}式可得
\begin{align*}
\|u\|_{L^q(\mathbb{R}^n)}
\leqslant \|u\|_{L^n(\mathbb{R}^n)}^\theta \|u\|_{L^{\frac{n^2}{n-1}}(\mathbb{R}^n)}^{1-\theta}
\leqslant C\|u\|_{W^{1,n}(\mathbb{R}^n)}.
\end{align*}
因此对一切 $q \in \left[n, \frac{n^2}{n-1}\right]$，都有
\begin{align}\label{eq::brezis-limit-interpolation}
\|u\|_{L^q(\mathbb{R}^n)} \leqslant C\|u\|_{W^{1,n}(\mathbb{R}^n)}.
\end{align}

现在用迭代法将 $q$ 的范围扩展至 $[n, +\infty)$。假设对某个整数 $m \geqslant n$，已经证明对一切 $q \in \left[n, \frac{mn}{n-1}\right]$ 有
\begin{align*}
\|u\|_{L^q(\mathbb{R}^n)} \leqslant C_m\|u\|_{W^{1,n}(\mathbb{R}^n)}.
\end{align*}
在 \eqref{eq::brezis-limit-young} 式中将 $m$ 换成 $m+1$，得到
\begin{align*}
\|u\|_{L^{\frac{(m+1)n}{n-1}}(\mathbb{R}^n)} \leqslant C_{m+1}\left(\|u\|_{L^{\frac{mn}{n-1}}(\mathbb{R}^n)} + \|Du\|_{L^n(\mathbb{R}^n)}\right) \leqslant C_{m+1}'\|u\|_{W^{1,n}(\mathbb{R}^n)}.
\end{align*}
再由\hyperref[theorem:广义Holder不等式]{广义Hölder不等式}，对一切 $q \in \left[n, \frac{(m+1)n}{n-1}\right]$，同理可得
\begin{align*}
\|u\|_{L^q(\mathbb{R}^n)} \leqslant C\|u\|_{W^{1,n}(\mathbb{R}^n)}.
\end{align*}
由归纳法，对任意给定的 $q \in [n, +\infty)$，取 $m$ 充分大使得 $\frac{mn}{n-1} \geqslant q$，即得
\begin{align}\label{eq::brezis-limit-final}
\|u\|_{L^q(\mathbb{R}^n)} \leqslant C\|u\|_{W^{1,n}(\mathbb{R}^n)}.
\end{align}

最后，由\hyperref[theorem:ch05-光滑函数的局部逼近定理]{光滑函数的局部逼近定理}，对任意 $u \in W^{1,n}(\mathbb{R}^n)$，存在 $u_k \in C_c^1(\mathbb{R}^n)$ 使得 $u_k \to u$ 于 $W^{1,n}(\mathbb{R}^n)$.并且
\begin{align}\label{eqr::Fjioj32334rrerr}
\|u_k\|_{L^q(\mathbb{R}^n)} \leqslant C\|u_k\|_{W^{1,n}(\mathbb{R}^n)}.
\end{align}
因为 $\{u_k\}$ 在 $W^{1,n}(\mathbb{R}^n)$ 中收敛，故为 Cauchy 列。由不等式\eqref{eqr::Fjioj32334rrerr}
可知，对任意 $k,j \geqslant 1$，有
\begin{align*}
\|u_k - u_j\|_{L^q(\mathbb{R}^n)} \leqslant C\|u_k - u_j\|_{W^{1,n}(\mathbb{R}^n)}.
\end{align*}
因此 $\{u_k\}$ 在 $L^q(\mathbb{R}^n)$ 中也为 Cauchy 列。由于 $L^q(\mathbb{R}^n)$ 完备，存在 $v \in L^q(\mathbb{R}^n)$ 使得
\begin{align*}
u_k \to v \quad \text{于} \ L^q(\mathbb{R}^n).
\end{align*}
另一方面，由 $u_k \to u$ 于 $W^{1,n}(\mathbb{R}^n)$ 可得 $u_k \to u$ 于 $L^n(\mathbb{R}^n)$。注意到 $q \geqslant n$，且所有 $u_k$ 的支集均包含在某个固定紧集内，因此由\cref{Real Analysis-proposition:Lp空间包含关系}知$u_k \to v$ 于 $L^q(\mathbb{R}^n)$ 蕴含 $u_k \to v$ 于 $L^n(\mathbb{R}^n)$。由 $L^n(\mathbb{R}^n)$ 中极限的唯一性，得到
\begin{align*}
v = u \quad \text{a.e.} \ \text{于} \ \mathbb{R}^n.
\end{align*}
因此 $u \in L^q(\mathbb{R}^n)$。在不等式
\begin{align*}
\|u_k\|_{L^q(\mathbb{R}^n)} \leqslant C\|u_k\|_{W^{1,n}(\mathbb{R}^n)}
\end{align*}
中令 $k \to \infty$，由 $L^q$ 范数与 $W^{1,n}$ 范数的连续性，即得
\begin{align*}
\|u\|_{L^q(\mathbb{R}^n)} \leqslant C\|u\|_{W^{1,n}(\mathbb{R}^n)}.
\end{align*}
\end{proof}

\begin{theorem}[$W^{1,p}$ 的估计，$1\leqslant p<n$]\label{theorem:ch05-W1p估计}
设 $U$ 为 $\mathbb{R}^{n}$ 的有界开子集，$\partial U$ 为 $C^{1}$。设 $1\leqslant p<n$，$u\in W^{1,p}(U)$。则 $u\in L^{p^{*}}(U)$，并有估计
\begin{align*}
\|u\|_{L^{p^{*}}(U)}\leqslant C\|u\|_{W^{1,p}(U)},
\end{align*}
其中常数 $C$ 只依赖于 $p$、$n$ 与 $U$。
\end{theorem}
\begin{proof}
由于 $\partial U$ 为 $C^{1}$，根据\hyperref[theorem:ch05-扩张定理]{扩张定理}存在延拓 $Eu=\overline{u}\in W^{1,p}(\mathbb{R}^{n})$，使
\begin{align}\label{eq:iwejofiw232333}
\begin{cases}
\text{在}\ U\ \text{中}\overline{u}=u,\overline{u}\ \text{有紧支集},\\
\|\overline{u}\|_{W^{1,p}(\mathbb{R}^{n})}\leqslant C\|u\|_{W^{1,p}(U)}.
\end{cases}
\end{align}
由于 $\overline{u}$ 有紧支集，由\hyperref[theorem:ch05-光滑函数的局部逼近定理]{光滑函数的局部逼近定理}存在函数 $u_{m}\in C^{\infty}_{c}(\mathbb{R}^{n})$（$m=1,2,\cdots$）使
\begin{align*}
u_{m}\to\overline{u}\quad\text{于}\ W^{1,p}(\mathbb{R}^{n}).
\end{align*}
现由\hyperref[theorem:Gagliardo--Nirenberg--Sobolev不等式]{Gagliardo--Nirenberg--Sobolev不等式}，对一切 $l,m\geqslant 1$ 有 $\|u_{m}-u_{l}\|_{L^{p^{*}}(\mathbb{R}^{n})}\leqslant C\|Du_{m}-Du_{l}\|_{L^{p}(\mathbb{R}^{n})}$。再结合上式可得
\begin{align*}
u_{m}\to\overline{u}\quad\text{于}\ L^{p^{*}}(\mathbb{R}^{n}).
\end{align*}
由于\hyperref[theorem:Gagliardo--Nirenberg--Sobolev不等式]{Gagliardo--Nirenberg--Sobolev不等式}还给出 $\|u_{m}\|_{L^{p^{*}}(\mathbb{R}^{n})}\leqslant C\|Du_{m}\|_{L^{p}(\mathbb{R}^{n})}$，上述两式给出界
\begin{align*}
\|\overline{u}\|_{L^{p^{*}}(\mathbb{R}^{n})}\leqslant C\|D\overline{u}\|_{L^{p}(\mathbb{R}^{n})}.
\end{align*}
此不等式与\eqref{eq:iwejofiw232333}式合起来即得
\begin{align*}
\|\overline{u}\|_{L^{p^{*}}(\mathbb{R}^{n})}\leqslant C\|D\overline{u}\|_{L^{p}(\mathbb{R}^{n})}\leqslant C\|\overline{u}\|_{W^{1,p}(\mathbb{R}^{n})}\leqslant C^2\|u\|_{W^{1,p}(U)}.
\end{align*}
\end{proof}

\begin{theorem}[$W^{1,p}_{0}$ 的估计，$1\leqslant p<n$]\label{theorem:ch05-W1p0估计}
设 $U$ 为 $\mathbb{R}^{n}$ 的有界开子集,$u\in W^{1,p}_{0}(U)$，其中 $1\leqslant p<n$。则对每个 $q\in[1,p^{*}]$ 有估计
\begin{align*}
\|u\|_{L^{q}(U)}\leqslant C\|Du\|_{L^{p}(U)},
\end{align*}
其中常数 $C$ 只依赖于 $p$、$q$、$n$ 与 $U$。
\end{theorem}
\begin{remark}
特别地,对于所有的 $q=p\in \left[1,\frac{n}{2}\right]$,因 $p<p^*$,由\cref{theorem:ch05-W1p0估计}得
\begin{align*}
\|u\|_{L^q(U)} \leqslant C\|Du\|_{L^q(U)}, \quad u \in W_0^{1,q}(U);
\end{align*}
而当 $q>\frac{n}{2}$ 时,取 $p=\frac{nq}{n+q}$,则 $q=p^*$,$\frac{1}{p}-\frac{1}{q}=\frac{1}{n}$,根据\hyperref[Real Analysis-theorem:实变函数--Hölder不等式]{Hölder不等式},若 $u\in W_0^{1,q}$,则有
\begin{align*}
\|Du\|_{L^p(U)} \leqslant |U|^{\frac{1}{n}}\|Du\|_{L^q(U)},
\end{align*}
再由\cref{theorem:ch05-W1p0估计}就得到 $n\geqslant 2$ 时
\begin{align*}
\|u\|_{L^q(U)} \leqslant C\|Du\|_{L^p(U)} \leqslant C|U|^{\frac{1}{n}}\|Du\|_{L^q(U)}, \quad u \in W_0^{1,q}(U), q \in \left(\frac{n}{2}, \infty \right).
\end{align*}
(这种类型的不等式称为 Poincaré 不等式,其他类型的 Poincaré 不等式将在后面的 §5.8.1 中建立.)它与\cref{theorem:ch05-W1p估计}的区别是不等式的右边仅出现 $u$ 的梯度. 但是不等式中的常数 $C$ 和\hyperref[theorem:Gagliardo--Nirenberg--Sobolev不等式]{Gagliardo--Nirenberg--Sobolev不等式}、\cref{theorem:ch05-W1p估计}、\cref{theorem:ch05-W1p0估计}中的一样,当 $p \to n^-$ 时,即 $q \to \infty$ 时是无界的.

而实际上 Poincaré 不等式中的常数 $C$ 可取为 $U$ 的半径,参见附录 §B.2.m 公式 (20).

对于有界的 $U$,由 Poincaré 不等式,$W_0^{1,p}(U)$ 上范数 $\|Du\|_{L^p(U)}$ 等价于范数 $\|u\|_{W^{1,p}(U)}$.
\end{remark}
\begin{proof}
由于 $u\in W^{1,p}_{0}(U)$，存在函数 $u_{m}\in C^{\infty}_{c}(U)$（$m=1,2,\cdots$）在 $W^{1,p}(U)$ 中收敛于 $u$。我们把每个 $u_{m}$ 在 $\mathbb{R}^{n}\setminus\overline{U}$ 上作零延拓，并应用\hyperref[theorem:Gagliardo--Nirenberg--Sobolev不等式]{Gagliardo--Nirenberg--Sobolev不等式}得
\begin{align}\label{eq::jjswjwiijjjjjjjj}
\left\| u_m \right\| _{L^{p^*}\left( U \right)}=\left\| u_m \right\| _{L^{p^*}\left( \mathbb{R} ^n \right)}\leqslant C\left\| Du_m \right\| _{L^p\left( \mathbb{R} ^n \right)}=C\left\| Du_m \right\| _{L^p\left( U \right)}.
\end{align}
因为 $u_m\to u$ 于 $W^{1,p}(U)$，零延拓后仍有 $u_m\to u$ 于 $W^{1,p}(\mathbb{R}^n)$，所以 $Du_m\to Du$ 于 $L^{p}(\mathbb{R}^n)$。另一方面，对任意 $l,m\geqslant 1$，由\hyperref[theorem:Gagliardo--Nirenberg--Sobolev不等式]{Gagliardo--Nirenberg--Sobolev不等式}有
\begin{align*}
\|u_m-u_l\|_{L^{p^{*}}(\mathbb{R}^n)}\leqslant C\|Du_m-Du_l\|_{L^{p}(\mathbb{R}^n)}.
\end{align*}
由于 $\{Du_m\}$ 在 $L^{p}(\mathbb{R}^n)$ 中收敛，故 $\{u_m\}$ 在 $L^{p^{*}}(\mathbb{R}^n)$ 中为 Cauchy 列，
由于 $L^{p^*}(\mathbb{R}^n)$ 完备，存在 $v\in L^{p^*}(\mathbb{R}^n)$ 使得
\begin{align*}
u_m\to v \quad \text{于 } L^{p^*}(\mathbb{R}^n).
\end{align*}
因为 $U$ 有界，$|U|<\infty$。对任意 $f\in L^{p^*}(U)$，由 \hyperref[Real Analysis-theorem:实变函数--Hölder不等式]{Hölder不等式}有
\begin{align*}
\|f\|_{L^p(U)}\leqslant |U|^{\frac{1}{p}-\frac{1}{p^*}}\|f\|_{L^{p^*}(U)}.
\end{align*}
因此 $u_m\to v$ 于 $L^p(U)$，零延拓后 $u_m\to v$ 于 $L^p(\mathbb{R}^n)$。另一方面，由 $u_m\to u$ 于 $W^{1,p}(U)$ 可得 $u_m\to u$ 于 $L^p(U)$，零延拓后 $u_m\to u$ 于 $L^p(\mathbb{R}^n)$。由 $L^p$ 极限的唯一性，$v=u$ 几乎处处于 $\mathbb{R}^n$。于是 $u\in L^{p^*}(U)$，且
\begin{align*}
u_m\to u \quad \text{于 } L^{p^*}(\mathbb{R}^n).
\end{align*}
于是 $u_m\to u$ 于 $L^{p^{*}}(\mathbb{R}^n)$。在不等式\eqref{eq::jjswjwiijjjjjjjj}中令 $m\to\infty$，即得
\begin{align*}
\|u\|_{L^{p^{*}}(U)}\leqslant C\|Du\|_{L^{p}(U)}.
\end{align*}
由于 $|U|<\infty$，当 $1\leqslant q\leqslant p^{*}$ 时,根据\hyperref[Real Analysis-theorem:实变函数--Hölder不等式]{Hölder不等式}还有
\begin{align*}
\|u\|_{L^{q}(U)}\leqslant |U|^{\frac{1}{q}-\frac{1}{p^{*}}}\|u\|_{L^{p^{*}}(U)}\leqslant C\|u\|_{L^{p^{*}}(U)}.
\end{align*}
\end{proof}
\begin{remark}
\textbf{边界情形 $p=n$.} 下面假设
\begin{align*}
p=n.
\end{align*}
由\cref{theorem:ch05-W1p估计} 及 $p^{*}=\dfrac{np}{n-p}\to+\infty$（当 $p\to n$）这一事实，我们或许会期望：若 $u\in W^{1,n}(U)$，则由\cref{theorem:ch05-W1p0估计}得$u\in L^{\infty}(U)$。然而当 $n>1$ 时这是错的：例如若 $U=B^{0}(0,1)$，函数 $u=\ln\ln\left(1+\dfrac{1}{|x|}\right)\in W^{1,n}(U)$,但不属于 $L^{\infty}(U)$。我们将在\cref{subsection:ch05-边界情形}中继续讨论这种边界情形。
\end{remark}


\subsection{Morrey 不等式}\label{subsection:ch05-Morrey}

我们将证明：若 $u\in W^{1,p}(U)$，则在一个零测集上重新定义之后，$u$ 事实上是H\"older连续的。

\begin{theorem}[Morrey不等式]\label{theorem:Morrey不等式}
设 $n<p\leqslant\infty$。则存在只依赖于 $p$ 与 $n$的常数$C$，使对一切 $u\in C^{1}(\mathbb{R}^{n})$ 有
\begin{align}\label{eq::chapter05jipde}
\|u\|_{C^{0,\gamma}(\mathbb{R}^{n})}\leqslant C\|u\|_{W^{1,p}(\mathbb{R}^{n})},
\end{align}
其中$\gamma\triangleq 1-\frac{n}{p}.$
\end{theorem}
\begin{proof}

1. 我们断言存在常数 $C$，只依赖于 $n$，使对每个球 $B(x,r)\subset\mathbb{R}^{n}$ 有
\begin{align}\label{eqwLJKjioefjioejw2io3}
\fint_{B(x,r)}|u(y)-u(x)|\,\mathrm{d}y
\leqslant \frac{1}{n\alpha(n)}\int_{B(x,r)}\frac{|Du(y)|}{|y-x|^{n-1}}\,\mathrm{d}y.
\end{align}
为证此，固定任一点 $w\in\partial B(0,1)$。若 $0<s<r$，则由微积分基本定理得
\begin{align*}
|u(x+sw)-u(x)|
=\left|\int_{0}^{s}\frac{\mathrm{d}}{\mathrm{d}t}u(x+tw)\,\mathrm{d}t\right|
=\left|\int_{0}^{s}w\cdot Du(x+tw)\,\mathrm{d}t\right|
\leqslant\int_{0}^{s}|Du(x+tw)|\,\mathrm{d}t.
\end{align*}
因而再由\hyperref[Real Analysis-theorem:实变函数-Fubini定理]{Fubini定理}可得
\begin{align*}
\int_{\partial B(0,1)}|u(x+sw)-u(x)|\,\mathrm{d}S(w)
\leqslant \int_{\partial B(0,1)}{\int_0^s{|Du(x+tw)|\,\mathrm{d}t}\,\mathrm{d}S(w)}=\int_{0}^{s}\int_{\partial B(0,1)}|Du(x+tw)|\,\mathrm{d}S(w)\,\mathrm{d}t.
\end{align*}
现取$y=x+tw$,则$t=|x-y|$,故
\begin{align*}
\int_{0}^{s}\int_{\partial B(0,1)}|Du(x+tw)|\,\mathrm{d}S(w)\,\mathrm{d}t
&=\int_{0}^{s}\int_{\partial B(x,t)}\frac{|Du(y)|}{t^{n-1}}\,\mathrm{d}S(y)\,\mathrm{d}t\\
&=\int_{B(x,s)}\frac{|Du(y)|}{|x-y|^{n-1}}\,\mathrm{d}y
\leqslant\int_{B(x,r)}\frac{|Du(y)|}{|x-y|^{n-1}}\,\mathrm{d}y.
\end{align*}
进一步取$z=x+sw$,则
\begin{align*}
\int_{\partial B(0,1)}|u(x+sw)-u(x)|\,\mathrm{d}S(w)
=\frac{1}{s^{n-1}}\int_{\partial B(x,s)}|u(z)-u(x)|\,\mathrm{d}S(z),
\end{align*}
把前两个计算代入上式，得估计
\begin{align*}
\int_{\partial B(x,s)}|u(z)-u(x)|\,\mathrm{d}S(z)
\leqslant s^{n-1}\int_{B(x,r)}\frac{|Du(y)|}{|x-y|^{n-1}}\,\mathrm{d}y.
\end{align*}
现对 $s$ 从 $0$ 到 $r$ 积分得
\begin{align*}
\int_{B(x,r)}|u(y)-u(x)|\,\mathrm{d}y
\leqslant\frac{r^{n}}{n}\int_{B(x,r)}\frac{|Du(y)|}{|x-y|^{n-1}}\,\mathrm{d}y.
\end{align*}
上式蕴含了\eqref{eqwLJKjioefjioejw2io3}式.

2. 现固定 $x\in\mathbb{R}^{n}$.由\hyperref[Real Analysis-theorem:实变函数--Hölder不等式]{Hölder不等式}得
\begin{align*}
\int_{B(x,r)} |u(y)|\,\mathrm{d}y
&\leqslant \left( \int_{B(x,r)} |u(y)|^p \,\mathrm{d}y \right)^{\frac{1}{p}} \left( \int_{B(x,r)} 1^{\frac{p}{p-1}} \,\mathrm{d}y \right)^{\frac{p-1}{p}} \\
&= \|u\|_{L^p(B(x,r))} |B(x,r)|^{\frac{p-1}{p}} = \|u\|_{L^p(B(x,r))} (\alpha(n)r^n)^{\frac{p-1}{p}}.
\end{align*}
因此
\begin{align}\label{eq::jijjswiiowioeioi3329iwiswiosmfc-1}
\fint_{B(x,r)} |u(y)|\,\mathrm{d}y
= \frac{1}{|B(x,r)|} \int_{B(x,r)} |u(y)|\,\mathrm{d}y
\leqslant \frac{1}{\alpha(n)r^n} \|u\|_{L^p(B(x,r))} (\alpha(n)r^n)^{\frac{p-1}{p}}
= \frac{r^{-\frac{n}{p}}}{\alpha^{\frac{1}{p}}(n)} \|u\|_{L^p(B(x,r))}.
\end{align}
由\hyperref[Real Analysis-theorem:实变函数--Hölder不等式]{Hölder不等式}得
\begin{align}\label{eqrLJIOJIOIOIO2903993}
\int_{B(x,r)} \frac{|Du(y)|}{|x-y|^{n-1}}\,\mathrm{d}y
\leqslant \left( \int_{B(x,r)} |Du(y)|^p \,\mathrm{d}y \right)^{\frac{1}{p}} \left( \int_{B(x,r)} |x-y|^{-(n-1)\frac{p}{p-1}} \,\mathrm{d}y \right)^{\frac{p-1}{p}}.
\end{align}
下面计算右端第二个积分. 令 $y=x+z$, 并作球坐标换元 $z=\rho\omega$, 其中 $\rho\in(0,r)$, $\omega\in\partial B(0,1)$. 则 $\mathrm{d}z=\rho^{n-1}\,\mathrm{d}\sigma(\omega)\,\mathrm{d}\rho$, 且 $|x-y|=|z|=\rho$. 于是
\begin{align*}
\int_{B(x,r)} |x-y|^{-(n-1)\frac{p}{p-1}} \,\mathrm{d}y
&= \int_{B(0,r)} |z|^{-(n-1)\frac{p}{p-1}} \,\mathrm{d}z = \int_{0}^{r} \int_{\partial B(0,1)} \rho^{-(n-1)\frac{p}{p-1}} \rho^{n-1} \,\mathrm{d}\sigma(\omega)\,\mathrm{d}\rho \\
&= n\alpha(n) \int_{0}^{r} \rho^{n-1-(n-1)\frac{p}{p-1}} \,\mathrm{d}\rho = n\alpha(n) \int_{0}^{r} \rho^{-\frac{n-1}{p-1}} \,\mathrm{d}\rho \\
&= n\alpha(n) \frac{p-1}{p-n} r^{\frac{p-n}{p-1}}.
\end{align*}
这里用到了 $p>n$ 以保证积分收敛. 将此结果代入不等式\eqref{eqrLJIOJIOIOIO2903993}中, 得到
\begin{align}
\frac{1}{n\alpha(n)}\int_{B(x,r)} \frac{|Du(y)|}{|x-y|^{n-1}}\,\mathrm{d}y
&\leqslant \frac{1}{n\alpha(n)}\|Du\|_{L^p(B(x,r))} \left( n\alpha(n) \frac{p-1}{p-n} r^{\frac{p-n}{p-1}} \right)^{\frac{p-1}{p}} \nonumber \\
&= \frac{\left(\frac{p-1}{p-n}\right)^{1-\frac{1}{p}} r^{1-\frac{n}{p}}}{(n\alpha(n))^{\frac{1}{p}}} \|Du\|_{L^p(B(x,r))}.\label{eq::jijjswiiowioeioi3329iwiswiosmfc-2}
\end{align}
再如下应用不等式\eqref{eqwLJKjioefjioejw2io3}及不等式\eqref{eq::jijjswiiowioeioi3329iwiswiosmfc-1}\eqref{eq::jijjswiiowioeioi3329iwiswiosmfc-2}, 得
\begin{align*}
|u(x)| &\leqslant \fint_{B(x,r)} |u(x)-u(y)|\,\mathrm{d}y + \fint_{B(x,r)} |u(y)|\,\mathrm{d}y \\
&\leqslant \frac{1}{n\alpha(n)} \int_{B(x,r)} \frac{|Du(y)|}{|x-y|^{n-1}}\,\mathrm{d}y + \frac{r^{-\frac{n}{p}}}{\alpha^{\frac{1}{p}}(n)} \|u\|_{L^p(B(x,r))} \\
&\leqslant \frac{\left(\frac{p-1}{p-n}\right)^{1-\frac{1}{p}} r^{1-\frac{n}{p}}}{(n\alpha(n))^{\frac{1}{p}}} \left(\int_{B(x,r)} |Du|^p\,\mathrm{d}y\right)^{\frac{1}{p}} + \frac{r^{-\frac{n}{p}}}{\alpha^{\frac{1}{p}}(n)} \|u\|_{L^p(B(x,r))} \\
&\triangleq C_1 r^{1-\frac{n}{p}} \|Du\|_{L^p(B(x,r))} + C_2 r^{-\frac{n}{p}} \|u\|_{L^p(B(x,r))}.
\end{align*}
因为 $x\in\mathbb{R}^n$ 是任意的, 可知
\begin{align}
\sup\limits_{\mathbb{R}^n} |u| \leqslant C_1 r^{1-\frac{n}{p}} \|Du\|_{L^p(\mathbb{R}^n)} + C_2 r^{-\frac{n}{p}} \|u\|_{L^p(\mathbb{R}^n)}. \label{eq:ch05-morrey-23}
\end{align}
再任选两点 $x,y\in\mathbb{R}^n$, 记 $r \triangleq |x-y|$. 令 $W \triangleq B(x,r)\cap B(y,r)$. 则
\begin{align}
|u(x)-u(y)| \leqslant \fint_W |u(x)-u(z)|\,\mathrm{d}z + \fint_W |u(y)-u(z)|\,\mathrm{d}z. \label{eq:ch05-morrey-24}
\end{align}
由\eqref{eqwLJKjioefjioejw2io3}及\hyperref[Real Analysis-theorem:实变函数--Hölder不等式]{Hölder不等式}, 类似可以得到
\begin{align}
\fint_W |u(x)-u(z)|\,\mathrm{d}z \leqslant \fint_{B(x,r)} |u(x)-u(z)|\,\mathrm{d}z \leqslant C_1 r^{1-\frac{n}{p}} \|Du\|_{L^p(B(x,r))}, \label{eq:ch05-morrey-25}
\end{align}
及
\begin{align*}
\fint_W |u(y)-u(z)|\,\mathrm{d}z \leqslant C_1 r^{1-\frac{n}{p}} \|Du\|_{L^p(B(y,r))}.
\end{align*}
同理, 注意到 $B(y,r)\subseteq B(x,2r)$, 我们得到
\begin{align*}
\fint_W |u(y)-u(z)|\,\mathrm{d}z \leqslant C_1 r^{1-\frac{n}{p}} \|Du\|_{L^p(B(x,2r))}.
\end{align*}
我们把此估计式以及 \eqref{eq:ch05-morrey-25} 式代入 \eqref{eq:ch05-morrey-24} 中得
\begin{align}\label{eq::jiwj333j33333}
|u(x)-u(y)| \leqslant 2C_1 r^{1-\frac{n}{p}} \|Du\|_{L^p(B(x,2r))}.
\end{align}
因此
\begin{align*}
[u]_{C^{0,1-\frac{n}{p}}(\mathbb{R}^n)} = \sup\limits_{x\neq y} \left\{ \frac{|u(x)-u(y)|}{|x-y|^{1-\frac{n}{p}}} \right\} \leqslant 2C_1 \|Du\|_{L^p(\mathbb{R}^n)}.
\end{align*}
从而, 由此不等式及 \eqref{eq:ch05-morrey-23} 式, 得
\begin{align*}
\|u\|_{C^{0,1-\frac{n}{p}}(\mathbb{R}^n)} \leqslant C_1 (r^{1-\frac{n}{p}}+2) \|Du\|_{L^p(\mathbb{R}^n)} + C_2 r^{-\frac{n}{p}} \|u\|_{L^p(\mathbb{R}^n)}.
\end{align*}
易知存在唯一的只依赖于 $n$ 与 $p$ 的 $r=r(n,p)>0$ 满足 $C_1 (r^{1-\frac{n}{p}}+2) = C_2 r^{-\frac{n}{p}}$. 故得
\begin{align*}
\|u\|_{C^{0,1-\frac{n}{p}}(\mathbb{R}^n)} \leqslant C_1 (r^{1-\frac{n}{p}}(n,p)+2) \|u\|_{W^{1,p}(\mathbb{R}^n)}.
\end{align*}
即得\eqref{eq::chapter05jipde}式, 其中 $C = C_1 (r^{1-\frac{n}{p}}(n,p)+2)$ 只依赖于 $n$ 与 $p$. (注意当 $p\to n$ 时, $C_1$ 是无界的.)

当 $p = \infty$ 时，$\gamma = 1 - \frac{n}{p} = 1$。此时不等式 \eqref{eq::chapter05jipde} 成为
\begin{align*}
\|u\|_{C^{0,1}(\mathbb{R}^n)} \leqslant C \|u\|_{W^{1,\infty}(\mathbb{R}^n)}.
\end{align*}
对任意 $x, y \in \mathbb{R}^n$，由微积分基本定理，有
\begin{align*}
u(y) - u(x) = \int_0^1 Du(x + t(y - x)) \cdot (y - x) \,\mathrm{d}t.
\end{align*}
从而
\begin{align*}
|u(y) - u(x)| \leqslant \int_0^1 |Du(x + t(y - x))| |y - x| \,\mathrm{d}t \leqslant \|Du\|_{L^\infty(\mathbb{R}^n)} |y - x|.
\end{align*}
因此
\begin{align*}
[u]_{C^{0,1}(\mathbb{R}^n)} = \sup_{x \neq y} \frac{|u(x) - u(y)|}{|x - y|} \leqslant \|Du\|_{L^\infty(\mathbb{R}^n)}.
\end{align*}
于是
\begin{align*}
\|u\|_{C^{0,1}(\mathbb{R}^n)} = \|u\|_{L^\infty(\mathbb{R}^n)} + [u]_{C^{0,1}(\mathbb{R}^n)} \leqslant \|u\|_{L^\infty(\mathbb{R}^n)} + \|Du\|_{L^\infty(\mathbb{R}^n)} \leqslant 2 \|u\|_{W^{1,\infty}(\mathbb{R}^n)}.
\end{align*}
取 $C = 2$，即得所需不等式。常数 $C$ 只依赖于 $n$，故满足定理要求。
\end{proof}

\begin{definition}\label{definition:ch05-版本}
如果$u=u^{*}$ a.e.,我们称 $u^{*}$ 为给定函数 $u$ 的一个\textbf{等价形}，
\end{definition}

\begin{theorem}[$W^{1,p}$ 的估计，$n<p\leqslant\infty$]\label{theorem:ch05-W1p有界估计}
设 $U$ 为 $\mathbb{R}^{n}$ 的有界开子集，$\partial U$ 为 $C^{1}$。设 $n<p\leqslant\infty$，$u\in W^{1,p}(U)$。则 $u$ 有一个等价形$u^{*}\in C^{0,\gamma}(\overline{U})$，其中 $\gamma=1-\dfrac{n}{p}$，并有估计
\begin{align*}
\|u^{*}\|_{C^{0,\gamma}(\overline{U})}\leqslant C\|u\|_{W^{1,p}(U)}.
\end{align*}
常数 $C$ 只依赖于 $p$、$n$ 与 $U$。
\end{theorem}
\begin{remark}
鉴于此定理，此后我们总把函数 $u\in W^{1,p}(U)$（$p>n$）等同于它的连续等价形。
\end{remark}
\begin{proof}
由于 $\partial U$ 为 $C^{1}$，根据\hyperref[theorem:ch05-扩张定理]{扩张定理}存在延拓 $Eu=\overline{u}\in W^{1,p}(\mathbb{R}^{n})$，使
\begin{align}
\begin{cases}
\overline{u}=u\ \text{于}\ U,\ \overline{u}\ \text{有紧支集},\\
\|\overline{u}\|_{W^{1,p}(\mathbb{R}^{n})}\leqslant C\|u\|_{W^{1,p}(U)}.
\end{cases}
\end{align}
先设 $n<p<\infty$。由于 $\overline{u}$ 有紧支集，由\hyperref[theorem:ch05-光滑函数的局部逼近定理]{光滑函数的局部逼近定理},取定理中的开集为 $U=\mathbb{R}^{n}$。此时边界 $\partial U=\varnothing$，故对任意 $\varepsilon>0$，$U_{\varepsilon}=\mathbb{R}^{n}$。设 $K\triangleq\operatorname{spt}\overline{u}$ 为紧集,则磨光函数
\begin{align*}
u_{\varepsilon}\triangleq\eta_{\varepsilon}*\overline{u}
\end{align*}
在整个 $\mathbb{R}^{n}$ 上有定义.由\hyperref[theorem:appC-光滑算子的性质]{光滑算子的性质\ref{theorem:appC-光滑算子的性质-6}}知$\operatorname{spt}u_{\varepsilon}\subseteq K+\overline{B(0,\varepsilon)}$，而 $K+\overline{B(0,\varepsilon)}$ 仍是紧集，所以 $u_{\varepsilon}\in C_{c}^{\infty}(\mathbb{R}^{n})$.又由
\begin{align*}
u_{\varepsilon}\to\overline{u}\quad\text{于}\ W^{1,p}_{\text{loc}}(\mathbb{R}^{n})
\end{align*}
可得
\begin{align*}
\left\| u_{\varepsilon}-\overline{u} \right\| _{W^{1,p}\left( \mathbb{R} ^n \right)}=\left\| u_{\varepsilon}-\overline{u} \right\| _{W^{1,p}\left( \mathrm{spt}\,u_{\varepsilon} \right)}\rightarrow 0,\quad \varepsilon \to 0.
\end{align*}
即存在函数 $u_{m}\in C^{\infty}_{c}(\mathbb{R}^{n})$ 使
\begin{align*}
u_{m}\to\overline{u}\quad\text{于}\ W^{1,p}(\mathbb{R}^{n}).
\end{align*}
现由\hyperref[theorem:Morrey不等式]{Morrey不等式}，对一切 $l,m\geqslant 1$ 有 $\|u_{m}-u_{l}\|_{C^{0,1-\frac{n}{p}}(\mathbb{R}^{n})}\leqslant C\|u_{m}-u_{l}\|_{W^{1,p}(\mathbb{R}^{n})}$，又\hyperref[theorem:ch05-holder空间Banach]{H\"older空间是Banach空间},故存在函数 $u^{*}\in C^{0,1-\frac{n}{p}}(\mathbb{R}^{n})$ 使
\begin{align*}
u_{m}\to u^{*}\quad\text{于}\ C^{0,1-\frac{n}{p}}(\mathbb{R}^{n}).
\end{align*}
因为 $u_{m}\to\overline{u}$ 于 $W^{1,p}(\mathbb{R}^{n})$，所以 $u_{m}\to\overline{u}$ 于 $L^{p}(\mathbb{R}^{n})$。因为 $u_{m}\to u^{*}$ 于 $C^{0,1-\frac{n}{p}}(\mathbb{R}^{n})$，而 $C^{0,1-\frac{n}{p}}$ 收敛蕴含一致收敛，进而蕴含 $L^{p}$ 收敛，所以 $u_{m}\to u^{*}$ 于 $L^{p}(\mathbb{R}^{n})$。由 $L^{p}$ 极限的唯一性，得到 $\overline{u}=u^{*}$ a.e. 于 $\mathbb{R}^{n}$。特别地，$u^{*}=u$ a.e. 于 $U$，故 $u^{*}$ 是 $u$ 的一个等价形。
由于\hyperref[theorem:Morrey不等式]{Morrey不等式}还给出 $\|u_{m}\|_{C^{0,1-\frac{n}{p}}(\mathbb{R}^{n})}\leqslant C\|u_{m}\|_{W^{1,p}(\mathbb{R}^{n})}$，不等式两边令 $m\to\infty$，得
\begin{align*}
\|u^{*}\|_{C^{0,1-\frac{n}{p}}(\mathbb{R}^{n})}\leqslant C\|\overline{u}\|_{W^{1,p}(\mathbb{R}^{n})}.
\end{align*}
此不等式与延拓估计合起来，即得
\begin{align*}
\left\| u^* \right\| _{C^{0,1-\frac{n}{p}}(\overline{U})}\leqslant \left\| u^* \right\| _{C^{0,1-\frac{n}{p}}\left( \mathbb{R} ^n \right)}\leqslant C\left\| \overline{u} \right\| _{W^{1,p}\left( \mathbb{R} ^n \right)}\leqslant C^2\left\| u \right\| _{W^{1,p}\left( U \right)},
\end{align*}
这就完成了 $n<p<\infty$ 情形的证明。

当\(p=\infty\)时,此时 \(\gamma=1-\frac{n}{p}=1\)，结论中的空间为 \(C^{0,1}(\overline U)\)，即 Lipschitz 连续函数空间。设 \(u\in W^{1,\infty}(U)\)。
由\hyperref[theorem:ch05-扩张定理]{扩张定理（\(p=\infty\) 情形）}，存在
\begin{align*}
\overline u=Eu\in W^{1,\infty}(\mathbb R^n),
\end{align*}
使得
\begin{align*}
\overline u=u\ \text{a.e. 于}\ U,\qquad \operatorname{spt}\overline u\ \text{紧},
\end{align*}
并且
\begin{align}\label{eq:W1inf-扩张估计}
\|\overline u\|_{W^{1,\infty}(\mathbb R^n)}
\leqslant C\|u\|_{W^{1,\infty}(U)}.
\end{align}

由\hyperref[theorem:Morrey不等式]{Morrey不等式（\(p=\infty\) 情形）}，\(\overline u\) 有一个 Lipschitz 连续等价形 \(u^*\in C^{0,1}(\mathbb R^n)\)，并且
\begin{align}\label{eq:W1inf-Morrey估计}
\|u^*\|_{C^{0,1}(\mathbb R^n)}
\leqslant C\|\overline u\|_{W^{1,\infty}(\mathbb R^n)}.
\end{align}
将 \(u^*\) 限制在 \(\overline U\) 上，得到 \(u^*|_{\overline U}\in C^{0,1}(\overline U)\)。由于 \(\overline u=u\) a.e. 于 \(U\)，且 \(u^*=\overline u\) a.e.，故
\begin{align*}
u^*=u\quad\text{a.e. 于}\ U,
\end{align*}
即 \(u^*\) 是 \(u\) 的连续等价形。
最后，由\eqref{eq:W1inf-扩张估计}\eqref{eq:W1inf-Morrey估计}式，得到
\begin{align*}
\|u^*\|_{C^{0,1}(\overline U)}
&\leqslant \|u^*\|_{C^{0,1}(\mathbb R^n)}
\leqslant C\|\overline u\|_{W^{1,\infty}(\mathbb R^n)}
\leqslant C^2\|u\|_{W^{1,\infty}(U)}.
\end{align*}
常数 \(C\) 只依赖于 \(n\) 与 \(U\)。这就完成了 \(p=\infty\) 情形的证明。
\end{proof}

\begin{corollary}\label{corollary:ch05-Morrey局部估计}
设 $n<p\leqslant\infty$，$x\in\mathbb{R}^n$，$r>0$。则存在只依赖于 $n$ 与 $p$ 的常数 $C$，使得对每个 $u\in W^{1,p}(B(x,2r))$，有
\begin{align}\label{eq:ch05-Morrey局部估计}
|u(y)-u(x)|\leqslant C r^{1-\frac{n}{p}}\|Du\|_{L^p(B(x,2r))}
\end{align}
对几乎所有的 $y\in B(x,r)$ 成立。特别地，若将右端的 $B(x,2r)$ 换为 $B(x,r)$，不等式仍然成立，即
\begin{align}\label{eq:ch05-Morrey局部估计-球内}
|u(y)-u(x)|\leqslant C r^{1-\frac{n}{p}}\|Du\|_{L^p(B(x,r))}
\end{align}
对几乎所有的 $y\in B(x,r)$ 成立。
\end{corollary}
\begin{proof}

1.先设 $u\in C^1(B(x,2r))$。由\hyperref[theorem:Morrey不等式]{Morrey不等式}证明中的估计式\eqref{eq::jiwj333j33333}，对任意 $y\in B(x,r)$，有
\begin{align*}
|u(y)-u(x)|\leqslant 2C_1 r^{1-\frac{n}{p}}\|Du\|_{L^p(B(x,2r))},
\end{align*}
其中 $C_1$ 只依赖于 $n$ 与 $p$。因此 \eqref{eq:ch05-Morrey局部估计}式在光滑情形下成立。

现设 $u\in W^{1,p}(B(x,2r))$。由\hyperref[theorem:Morrey不等式]{Morrey不等式}，存在常数 $C_M$ 使得对一切 $v\in W^{1,p}(B(x,2r))$，
\begin{align}\label{eq:ch05-Morrey-光滑}
\|v\|_{C^{0,\gamma}(B(x,2r))}\leqslant C_M\|v\|_{W^{1,p}(B(x,2r))},
\end{align}
其中 $\gamma=1-\frac{n}{p}$.由\hyperref[theorem:ch05-直到边界光滑的函数整体逼近定理]{边界光滑函数的整体逼近定理},可取磨光序列 $u_m\in C^\infty(\overline{B(x,2r)})$，使得
\begin{align}\label{eq::jioijiowio38ijoOPPA}
u_m\to u\quad\text{于}\ W^{1,p}(B(x,2r)).
\end{align}
由\cref{theorem:ch05-W1p有界估计}，取 $U=B(x,2r)$，则$u$ 有一个等价形 $u^*\in C^{0,\gamma}(\overline{U})$，并有估计
\begin{align}\label{eq::jioijiowio38ijoOPPA-1}
\|u_m-u^*\|_{C^{0,\gamma}(\overline{B(x,2r)})}
&\leqslant C_0\|u_m-u\|_{W^{1,p}(B(x,2r))},
\end{align}
其中常数 $C_0$ 只依赖于 $p$、$n$ 与 $U$（即只依赖于 $r$）。从而由\eqref{eq::jioijiowio38ijoOPPA}\eqref{eq::jioijiowio38ijoOPPA-1}式知$u_m\to u^*$ 在 $C^{0,\gamma}(\overline{B(x,2r)})$ 中成立。特别地，对每个 $y\in B(x,r)$，有
\begin{align*}
u_m(y)\to u^*(y),\qquad u_m(x)\to u^*(x).
\end{align*}
另一方面，因为 $u_m\to u$ 于 $W^{1,p}(B(x,2r))$，所以 $Du_m\to Du$ 于 $L^p(B(x,2r))$.再由 $L^p$ 范数的连续性得
\begin{align*}
\|Du_m\|_{L^p(B(x,2r))}\to \|Du\|_{L^p(B(x,2r))}.
\end{align*}
对每个 $m$，由光滑情形估计式\eqref{eq:ch05-Morrey-光滑}，得
\begin{align*}
|u_m(y)-u_m(x)|\leqslant C_1 r^{1-\frac{n}{p}}\|Du_m\|_{L^p(B(x,2r))}.
\end{align*}
令 $m\to\infty$，即得
\begin{align*}
|u^*(y)-u^*(x)|\leqslant C_1 r^{1-\frac{n}{p}}\|Du\|_{L^p(B(x,2r))}.
\end{align*}
由于 $u^*$ 是 $u$ 的等价形，故对几乎所有的 $y\in B(x,r)$，有 $u(y)=u^*(y)$ 且 $u(x)=u^*(x)$。因此 \eqref{eq:ch05-Morrey局部估计}式成立。

2.设 \(u\in W^{1,p}(B(0,1))\)。令
\begin{align*}
v\triangleq u-(u)_{B(0,1)},
\end{align*}
其中 \((u)_{B(0,1)}=\fint_{B(0,1)}u\,\mathrm{d}z\)。则 \((v)_{B(0,1)}=0\)。由\hyperref[theorem:ch05-Poincare]{Poincaré不等式}，存在只依赖于 \(n\) 与 \(p\) 的常数 \(C_P\)，使得
\begin{align}\label{eq:ch05-Morrey-Poincare-单位球}
\|v\|_{L^p(B(0,1))}\leqslant C_P\|Dv\|_{L^p(B(0,1))}.
\end{align}
由\cref{theorem:ch05-W1p有界估计}，取 \(U=B(0,1)\)，则存在只依赖于 \(n\) 与 \(p\) 的常数 \(C_M\)，以及 \(v\) 的连续等价形 \(v^*\in C^{0,\gamma}(\overline{B(0,1)})\)，使得
\begin{align}\label{eq:ch05-Morrey-嵌入-单位球-修正}
\|v^*\|_{C^{0,\gamma}(\overline{B(0,1)})}\leqslant C_M\|v\|_{W^{1,p}(B(0,1))},
\end{align}
其中 \(\gamma=1-\frac{n}{p}\)。由 \eqref{eq:ch05-Morrey-Poincare-单位球}得
\begin{align*}
\|v\|_{W^{1,p}(B(0,1))}
=\|v\|_{L^p(B(0,1))}+\|Dv\|_{L^p(B(0,1))}\leqslant (C_P+1)\|Dv\|_{L^p(B(0,1))}.
\end{align*}
代入 \eqref{eq:ch05-Morrey-嵌入-单位球-修正}，得
\begin{align*}
\|v^*\|_{C^{0,\gamma}(\overline{B(0,1)})}\leqslant C_M(C_P+1)\|Du\|_{L^p(B(0,1))}.
\end{align*}
因此，对任意 \(z\in B(0,1)\)，有
\begin{align*}
|u(z)-u(0)|=|v(z)-v(0)|=|v^*(z)-v^*(0)|
\leqslant \|v^*\|_{C^{0,\gamma}(\overline{B(0,1)})}|z|^\gamma\leqslant C_M(C_P+1)\|Du\|_{L^p(B(0,1))}.
\end{align*}
记 \(C_0\triangleq C_M(C_P+1)\)，则
\begin{align}\label{eq:ch05-Morrey-单位球估计-修正}
|u(z)-u(0)|\leqslant C_0\|Du\|_{L^p(B(0,1))},\quad \forall z\in B(0,1).
\end{align}

设 \(u\in W^{1,p}(B(x,r))\)。定义
\begin{align*}
w(z)\triangleq u(x+rz),\quad z\in B(0,1).
\end{align*}
则 \(w\in W^{1,p}(B(0,1))\)，且
\begin{align*}
Dw(z)=r\,Du(x+rz).
\end{align*}
计算 \(L^p\) 范数：
\begin{align*}
\|Dw\|_{L^p(B(0,1))}^p
=\int_{B(0,1)}|r\,Du(x+rz)|^p\,\mathrm{d}z=r^p\int_{B(0,1)}|Du(x+rz)|^p\,\mathrm{d}z.
\end{align*}
作变量替换 \(y=x+rz\)，则 \(\mathrm{d}z=r^{-n}\mathrm{d}y\)，积分区域变为 \(B(x,r)\)。于是
\begin{align*}
\|Dw\|_{L^p(B(0,1))}^p
=r^p\cdot r^{-n}\int_{B(x,r)}|Du(y)|^p\,\mathrm{d}y=r^{p-n}\|Du\|_{L^p(B(x,r))}^p.
\end{align*}
因此
\begin{align}\label{eq:ch05-Morrey-缩放范数-修正}
\|Dw\|_{L^p(B(0,1))}=r^{1-\frac{n}{p}}\|Du\|_{L^p(B(x,r))}.
\end{align}
对 \(w\) 应用单位球估计 \eqref{eq:ch05-Morrey-单位球估计-修正}，得到对任意 \(z\in B(0,1)\)，
\begin{align*}
|w(z)-w(0)|\leqslant C_0\|Dw\|_{L^p(B(0,1))}.
\end{align*}
令 \(y=x+rz\)，则 \(z=\frac{y-x}{r}\)，且 \(z\in B(0,1)\) 当且仅当 \(y\in B(x,r)\)。于是
\begin{align*}
|u(y)-u(x)|=|w(z)-w(0)|\leqslant C_0\|Dw\|_{L^p(B(0,1))}.
\end{align*}
代入 \eqref{eq:ch05-Morrey-缩放范数-修正}，即得
\begin{align*}
|u(y)-u(x)|\leqslant C_0 r^{1-\frac{n}{p}}\|Du\|_{L^p(B(x,r))}.
\end{align*}
取 \(C=C_0\)，便得到 \eqref{eq:ch05-Morrey局部估计-球内}。
\end{proof}


\subsection{一般的Sobolev不等式}\label{subsection:ch05-一般Sobolev不等式}

\begin{theorem}[一般的Sobolev不等式]\label{theorem:ch05-一般Sobolev不等式}设 $U$为$\mathbb{R}^{n}$ 的有界开子集，边界为 $C^{1}$，$u\in W^{k,p}(U)$。
\begin{enumerate}[(i)]
\item 若$k<\frac{n}{p},$则 $u\in L^{q}(U)$，其中$\frac{1}{q}=\frac{1}{p}-\frac{k}{n}.$
此外有估计
\begin{align}\label{eq::::--328393838382349892jsdjsu}
\|u\|_{L^{q}(U)}\leqslant C\|u\|_{W^{k,p}(U)},
\end{align}
常数 $C$ 只依赖于 $k$、$p$、$n$ 与 $U$。
\item 若$k>\frac{n}{p},$则 $u\in C^{k-\left[\frac{n}{p}\right]-1,\gamma}(\overline{U})$，其中
\begin{align*}
\gamma=
\begin{cases}
\left[\frac{n}{p}\right]+1-\dfrac{n}{p}, & \text{若 } \frac{n}{p}\ \text{不是整数},\\
\text{任一小于 }1\ \text{的正数}, & \text{若 } \frac{n}{p}\ \text{是整数}.
\end{cases}
\end{align*}
此外有估计
\begin{align*}
\|u\|_{C^{k-\left[\frac{n}{p}\right]-1,\gamma}(\overline{U})}\leqslant C\|u\|_{W^{k,p}(U)},
\end{align*}
常数 $C$ 只依赖于 $k$、$p$、$n$、$\gamma$ 与 $U$。
\end{enumerate}
\end{theorem}
\begin{proof}

1. 设 $k<\frac{n}{p}$。由于 $\partial U$ 为 $C^{1}$，由\hyperref[theorem:ch05-扩张定理]{扩张定理}，存在 $\overline{u}\in W^{k,p}(\mathbb{R}^{n})$，使得
\begin{align}\label{eq::jijwIIJEI9}
\overline{u}=u\ \text{于}\ U,\qquad \operatorname{supp}\overline{u}\ \text{紧},\qquad
\|\overline{u}\|_{W^{k,p}(\mathbb{R}^{n})}\leqslant C_{0}\|u\|_{W^{k,p}(U)}.
\end{align}
由\hyperref[theorem:ch05-光滑函数的局部逼近定理]{光滑函数的局部逼近定理}，取开集为 $\mathbb{R}^{n}$，存在 $u_{m}\in C_{c}^{\infty}(\mathbb{R}^{n})$，使得
\begin{align*}
u_{m}\to\overline{u}\quad\text{于}\ W^{k,p}(\mathbb{R}^{n}).
\end{align*}

先对 $u_{m}$ 证明相应估计。令
\begin{align*}
p_{0}\triangleq p,\qquad \frac{1}{p_{j+1}}\triangleq\frac{1}{p_{j}}-\frac{1}{n},\quad j=0,1,\cdots,k-1.
\end{align*}
由于 $k<\frac{n}{p}$，故对每个 $j=0,1,\cdots,k-1$，有
\begin{align*}
j+1\leqslant k<\frac{n}{p}\Longrightarrow \frac{1}{p}>\frac{j+1}{n}\Longrightarrow \frac{1}{p_j}=\frac{1}{p}-\frac{j}{n}>\frac{1}{n}>0,
\end{align*}
从而 $p_{j}<n$。因此\hyperref[theorem:Gagliardo--Nirenberg--Sobolev不等式]{Gagliardo--Nirenberg--Sobolev不等式}可以应用于指数 $p_{j}$。

对任意多重指标 $|\beta|\leqslant k-1$，令 $v\triangleq D^{\beta}u_{m}$。由于 $u_{m}\in C_{c}^{\infty}(\mathbb{R}^{n})$，故 $v\in C_{c}^{\infty}(\mathbb{R}^{n})$。对 $v$ 应用\hyperref[theorem:Gagliardo--Nirenberg--Sobolev不等式]{Gagliardo--Nirenberg--Sobolev不等式}，得到
\begin{align*}
\|D^{\beta}u_{m}\|_{L^{p_{1}}(\mathbb{R}^{n})}
=\|v\|_{L^{(p_{0})^{*}}(\mathbb{R}^{n})}\leqslant C_{1}\|Dv\|_{L^{p_{0}}(\mathbb{R}^{n})}=C_{1}\|D(D^{\beta}u_{m})\|_{L^{p}(\mathbb{R}^{n})}\leqslant C_{1}\|u_{m}\|_{W^{k,p}(\mathbb{R}^{n})}.
\end{align*}
因此 $u_{m}\in W^{k-1,p_{1}}(\mathbb{R}^{n})$，且
\begin{align*}
\left\| u_m \right\| _{W^{k-1,p_1}\left( \mathbb{R} ^n \right)}&=\left( \sum_{\left| \beta \right|\leqslant k-1}{\left\| D^{\beta}u_m \right\| _{L^{p_1}\left( \mathbb{R} ^n \right)}^{p_1}} \right) ^{\frac{1}{p_1}}\leqslant \left( \sum_{\left| \beta \right|\leqslant k-1}{C_1\left\| u_m \right\| _{W^{k,p}\left( \mathbb{R} ^n \right)}^{p_1}} \right) ^{\frac{1}{p_1}}
\\
&=C_{1}^{^{\frac{1}{p_1}}}\left( \sum_{\left| \beta \right|\leqslant k-1}{1} \right) ^{\frac{1}{p_1}}\left\| u_m \right\| _{W^{k,p}\left( \mathbb{R} ^n \right)}\triangleq C_{1}^{\prime}\left\| u_m \right\| _{W^{k,p}\left( \mathbb{R} ^n \right)}.
\end{align*}
假设对某个 $j\in\{0,1,\cdots,k-1\}$，已有
\begin{align*}
u_{m}\in W^{k-j,p_{j}}(\mathbb{R}^{n}),
\end{align*}
且
\begin{align*}
\|u_{m}\|_{W^{k-j,p_{j}}(\mathbb{R}^{n})}\leqslant C_{j}'\|u_{m}\|_{W^{k,p}(\mathbb{R}^{n})}.
\end{align*}
对任意多重指标 $|\beta|\leqslant k-j-1$，令 $v\triangleq D^{\beta}u_{m}$。因为 $u_{m}\in W^{k-j,p_{j}}(\mathbb{R}^{n})$，所以 $D^{\beta}u_{m}\in W^{1,p_{j}}(\mathbb{R}^{n})$，且 $v\in C_{c}^{\infty}(\mathbb{R}^{n})$。对 $v$ 应用\hyperref[theorem:Gagliardo--Nirenberg--Sobolev不等式]{Gagliardo--Nirenberg--Sobolev不等式}，得到
\begin{align*}
\|D^{\beta}u_{m}\|_{L^{p_{j+1}}(\mathbb{R}^{n})}
&=\|v\|_{L^{(p_{j})^{*}}(\mathbb{R}^{n})}\leqslant C_{j+1}\|Dv\|_{L^{p_{j}}(\mathbb{R}^{n})}=C_{j+1}\|D(D^{\beta}u_{m})\|_{L^{p_{j}}(\mathbb{R}^{n})}\\
&\leqslant C_{j+1}\|u_{m}\|_{W^{k-j,p_{j}}(\mathbb{R}^{n})}\leqslant C_{j+1}C_{j}'\|u_{m}\|_{W^{k,p}(\mathbb{R}^{n})}.
\end{align*}
因此 $u_{m}\in W^{k-j-1,p_{j+1}}(\mathbb{R}^{n})$，且
\begin{align*}
\left\| u_m \right\| _{W^{k-j-1,p_{j+1}}\left( \mathbb{R} ^n \right)}&=\left( \sum_{\left| \beta \right|\leqslant k-j-1}{\left\| D^{\beta}u_m \right\| _{L^{p_j+1}\left( \mathbb{R} ^n \right)}^{p_j+1}} \right) ^{\frac{1}{p_j+1}}\leqslant \left( \sum_{\left| \beta \right|\leqslant k-j-1}{C_{j+1}C_{j}^{\prime}\left\| u_m \right\| _{W^{k,p}\left( \mathbb{R} ^n \right)}^{p_j+1}} \right) ^{\frac{1}{p_j+1}}
\\
&=\left( C_{j+1}C_{j}^{\prime} \right) ^{\frac{1}{p_j+1}}\left( \sum_{\left| \beta \right|\leqslant k-j-1}{1} \right) ^{\frac{1}{p_j+1}}\left\| u_m \right\| _{W^{k,p}\left( \mathbb{R} ^n \right)}\triangleq C_{j+1}^{\prime}\left\| u_m \right\| _{W^{k,p}\left( \mathbb{R} ^n \right)}.
\end{align*}
重复上述论证，经过 $k$ 步后,即取$j=k-1$,得到
\begin{align*}
u_{m}\in W^{0,p_{k}}(\mathbb{R}^{n})=L^{p_{k}}(\mathbb{R}^{n})=L^{q}(\mathbb{R}^{n}),
\end{align*}
其中$\frac{1}{p_{k}}=\frac{1}{p}-\frac{k}{n}=\frac{1}{q}$.
并且每一步的估计逐次复合，得到
\begin{align*}
\left\| u_m \right\| _{L^q(\mathbb{R} ^n)}=\left\| u_m \right\| _{W^{0,p_k}\left( \mathbb{R} ^n \right)}\leqslant C_{k}^{\prime}\left\| u_m \right\| _{W^{k,p}(\mathbb{R} ^n)}\triangleq C\left\| u_m \right\| _{W^{k,p}(\mathbb{R} ^n)},
\end{align*}
其中常数 $C$ 只依赖于 $k$、$p$、$n$。

由于 $u_{m}\to\overline{u}$ 于 $W^{k,p}(\mathbb{R}^{n})$，故 $\{u_{m}\}$ 在 $W^{k,p}$ 中为 Cauchy 列。由上述估计可知，对任意 $l,m\geqslant1$，有
\begin{align*}
\|u_{m}-u_{l}\|_{L^{q}(\mathbb{R}^{n})}\leqslant C\|u_{m}-u_{l}\|_{W^{k,p}(\mathbb{R}^{n})}.
\end{align*}
因此 $\{u_{m}\}$ 在 $L^{q}$ 中也为 Cauchy 列。由于 $L^{q}(\mathbb{R}^{n})$ 完备，存在 $v\in L^{q}(\mathbb{R}^{n})$ 使得
\begin{align*}
u_{m}\to v\quad\text{于}\ L^{q}(\mathbb{R}^{n}).
\end{align*}
另一方面，由 $u_{m}\to\overline{u}$ 于 $W^{k,p}(\mathbb{R}^{n})$，可得 $u_{m}\to\overline{u}$ 于 $L^{p}(\mathbb{R}^{n})$。由$\frac{1}{q}=\frac{1}{p}-\frac{k}{n}<\frac{1}{p}$知$q>p$，且 $\overline{u}$ 具有紧支集，所以由\cref{Real Analysis-proposition:Lp空间包含关系}知$L^{q}$ 收敛蕴含 $L^{p}$ 收敛。由 $L^{p}$极限的唯一性，得到
\begin{align*}
v=\overline{u}\quad\text{a.e.}\,x\in \mathbb{R}^n.
\end{align*}
因此 $\overline{u}\in L^{q}(\mathbb{R}^{n})$，且在估计式中令 $m\to\infty$，得到
\begin{align*}
\|\overline{u}\|_{L^{q}(\mathbb{R}^{n})}\leqslant C\|\overline{u}\|_{W^{k,p}(\mathbb{R}^{n})}.
\end{align*}
再结合扩张估计\eqref{eq::jijwIIJEI9}，便得到
\begin{align*}
\|u\|_{L^{q}(U)}
\leqslant\|\overline{u}\|_{L^{q}(\mathbb{R}^{n})}\leqslant C\|\overline{u}\|_{W^{k,p}(\mathbb{R}^{n})}\leqslant C C_{0}\|u\|_{W^{k,p}(U)}.
\end{align*}
上述常数 $C C_{0}$ 由每一步的常数复合而成，只依赖于 $k$、$p$、$n$ 与 $U$。

2. 当(ii)中$p=\infty$时,结论就是\hyperref[theorem:Morrey不等式]{Morrey不等式}中$p=\infty$的情形,显然成立.故下设$p<\infty$.

现设条件 $k>\frac{n}{p}$ 成立且 $\frac{n}{p}$ 不是整数. 令
\begin{align*}
l\triangleq\left[\frac{n}{p}\right],
\end{align*}
则
\begin{align*}
l<\frac{n}{p}<l+1.
\end{align*}
由于 $\partial U$ 为 $C^{1}$, 由\hyperref[theorem:ch05-扩张定理]{扩张定理}, 存在 $\overline{u}\in W^{k,p}(\mathbb{R}^{n})$, 使得
\begin{align}\label{eq::jijwIIJEI9-2}
\overline{u}=u\ \text{于}\ U,\qquad \operatorname{supp}\overline{u}\ \text{紧},\qquad
\|\overline{u}\|_{W^{k,p}(\mathbb{R}^{n})}\leqslant C_{0}\|u\|_{W^{k,p}(U)}.
\end{align}
由\hyperref[theorem:ch05-光滑函数的局部逼近定理]{光滑函数的局部逼近定理}, 取开集为 $\mathbb{R}^{n}$, 存在 $u_{m}\in C_{c}^{\infty}(\mathbb{R}^{n})$, 使得
\begin{align}\label{eq::jijwIIJEI9-2-m}
u_{m}\to\overline{u}\quad\text{于}\ W^{k,p}(\mathbb{R}^{n}).
\end{align}
特别地, $\{u_{m}\}$ 在 $W^{k,p}(\mathbb{R}^{n})$ 中为 Cauchy 列, 且存在常数 $M>0$ 使得
\begin{align*}
\|u_{m}\|_{W^{k,p}(\mathbb{R}^{n})}\leqslant M,\quad m=1,2,\cdots.
\end{align*}

先对 $u_{m}$ 证明相应估计. 令
\begin{align*}
p_{0}\triangleq p,\qquad \frac{1}{p_{j+1}}\triangleq\frac{1}{p_{j}}-\frac{1}{n},\quad j=0,1,\cdots,l-1.
\end{align*}
由于 $l<\frac{n}{p}$, 故对每个 $j=0,1,\cdots,l-1$, 有
\begin{align*}
j+1\leqslant l<\frac{n}{p}\Longrightarrow \frac{1}{p}>\frac{j+1}{n}\Longrightarrow \frac{1}{p_j}=\frac{1}{p}-\frac{j}{n}>\frac{1}{n}>0,
\end{align*}
从而 $p_{j}<n$. 因此\hyperref[theorem:Gagliardo--Nirenberg--Sobolev不等式]{Gagliardo--Nirenberg--Sobolev不等式}可以应用于指数 $p_{j}$.

对任意多重指标 $|\beta|\leqslant k-1$, 令 $v\triangleq D^{\beta}u_{m}$. 由于 $u_{m}\in C_{c}^{\infty}(\mathbb{R}^{n})$, 故 $v\in C_{c}^{\infty}(\mathbb{R}^{n})$. 对 $v$ 应用\hyperref[theorem:Gagliardo--Nirenberg--Sobolev不等式]{Gagliardo--Nirenberg--Sobolev不等式}, 得到
\begin{align}\label{eq::jijwIIJEI9-2-first}
\|D^{\beta}u_{m}\|_{L^{p_{1}}(\mathbb{R}^{n})}
=\|v\|_{L^{(p_{0})^{*}}(\mathbb{R}^{n})}\leqslant C_{1}\|Dv\|_{L^{p_{0}}(\mathbb{R}^{n})}=C_{1}\|D(D^{\beta}u_{m})\|_{L^{p}(\mathbb{R}^{n})}\leqslant C_{1}\|u_{m}\|_{W^{k,p}(\mathbb{R}^{n})}.
\end{align}
因此 $u_{m}\in W^{k-1,p_{1}}(\mathbb{R}^{n})$, 且
\begin{align*}
\left\| u_m \right\| _{W^{k-1,p_1}\left( \mathbb{R} ^n \right)}&=\left( \sum_{\left| \beta \right|\leqslant k-1}{\left\| D^{\beta}u_m \right\| _{L^{p_1}\left( \mathbb{R} ^n \right)}^{p_1}} \right) ^{\frac{1}{p_1}}\leqslant \left( \sum_{\left| \beta \right|\leqslant k-1}{C_1^{p_1}\left\| u_m \right\| _{W^{k,p}\left( \mathbb{R} ^n \right)}^{p_1}} \right) ^{\frac{1}{p_1}}
\\&=C_{1}\left( \sum_{\left| \beta \right|\leqslant k-1}{1} \right) ^{\frac{1}{p_1}}\left\| u_m \right\| _{W^{k,p}\left( \mathbb{R} ^n \right)}\triangleq C_{1}^{\prime}\left\| u_m \right\| _{W^{k,p}\left( \mathbb{R} ^n \right)}.
\end{align*}
假设对某个 $j\in\{0,1,\cdots,l-1\}$, 已有
\begin{align*}
u_{m}\in W^{k-j,p_{j}}(\mathbb{R}^{n}),
\end{align*}
且
\begin{align*}
\|u_{m}\|_{W^{k-j,p_{j}}(\mathbb{R}^{n})}\leqslant C_{j}'\|u_{m}\|_{W^{k,p}(\mathbb{R}^{n})}.
\end{align*}
对任意多重指标 $|\beta|\leqslant k-j-1$, 令 $v\triangleq D^{\beta}u_{m}$. 因为 $u_{m}\in W^{k-j,p_{j}}(\mathbb{R}^{n})$, 所以 $D^{\beta}u_{m}\in W^{1,p_{j}}(\mathbb{R}^{n})$, 且 $v\in C_{c}^{\infty}(\mathbb{R}^{n})$. 对 $v$ 应用\hyperref[theorem:Gagliardo--Nirenberg--Sobolev不等式]{Gagliardo--Nirenberg--Sobolev不等式}, 得到
\begin{align*}
\|D^{\beta}u_{m}\|_{L^{p_{j+1}}(\mathbb{R}^{n})}
&=\|v\|_{L^{(p_{j})^{*}}(\mathbb{R}^{n})}\leqslant C_{j+1}\|Dv\|_{L^{p_{j}}(\mathbb{R}^{n})}=C_{j+1}\|D(D^{\beta}u_{m})\|_{L^{p_{j}}(\mathbb{R}^{n})}\\
&\leqslant C_{j+1}\|u_{m}\|_{W^{k-j,p_{j}}(\mathbb{R}^{n})}\leqslant C_{j+1}C_{j}'\|u_{m}\|_{W^{k,p}(\mathbb{R}^{n})}.
\end{align*}
因此 $u_{m}\in W^{k-j-1,p_{j+1}}(\mathbb{R}^{n})$, 且
\begin{align*}
\left\| u_m \right\| _{W^{k-j-1,p_{j+1}}\left( \mathbb{R} ^n \right)}
&=\left( \sum_{\left| \beta \right|\leqslant k-j-1}{\left\| D^{\beta}u_m \right\| _{L^{p_{j+1}}\left( \mathbb{R} ^n \right)}^{p_{j+1}}} \right) ^{\frac{1}{p_{j+1}}}\leqslant \left( \sum_{\left| \beta \right|\leqslant k-j-1}{\left( C_{j+1}C_{j}^{\prime}\left\| u_m \right\| _{W^{k,p}\left( \mathbb{R} ^n \right)} \right) ^{p_{j+1}}} \right) ^{\frac{1}{p_{j+1}}}\\
&=C_{j+1}C_{j}^{\prime}\left( \sum_{\left| \beta \right|\leqslant k-j-1}{1} \right) ^{\frac{1}{p_{j+1}}}\left\| u_m \right\| _{W^{k,p}\left( \mathbb{R} ^n \right)}\triangleq C_{j+1}^{\prime}\left\| u_m \right\| _{W^{k,p}\left( \mathbb{R} ^n \right)}.
\end{align*}
重复上述论证, 经过 $l$ 步后,即取$j=l-1$,得到
\begin{align*}
u_{m}\in W^{k-l,p_{l}}(\mathbb{R}^{n}),
\end{align*}
其中
\begin{align*}
\frac{1}{p_{l}}=\frac{1}{p}-\frac{l}{n}=\frac{1}{r},
\end{align*}
即 $p_{l}=r$, 且
\begin{align}\label{eq::jijwIIJEI9-2-l}
\|u_{m}\|_{W^{k-l,r}(\mathbb{R}^{n})}\leqslant C_{l}'\|u_{m}\|_{W^{k,p}(\mathbb{R}^{n})}.
\end{align}
因为 $l<\frac{n}{p}<l+1$, 所以
\begin{align*}
r=\frac{pn}{n-pl}=\frac{n}{\frac{n}{p}-l}>n.
\end{align*}

由 \eqref{eq::jijwIIJEI9-2-m}式知 $u_{m}\to\overline{u}$ 于 $W^{k,p}(\mathbb{R}^{n})$, 故 $\{u_{m}\}$ 在 $W^{k,p}(\mathbb{R}^{n})$ 中为 Cauchy 列. 由 \eqref{eq::jijwIIJEI9-2-l}式知, 对任意 $m,m'\geqslant1$ 有
\begin{align}\label{eq::jijwIIJEI9-2-cauchy}
\|u_{m}-u_{m'}\|_{W^{k-l,r}(\mathbb{R}^{n})}\leqslant C_{l}'\|u_{m}-u_{m'}\|_{W^{k,p}(\mathbb{R}^{n})}.
\end{align}
因此 $\{u_{m}\}$ 在 $W^{k-l,r}(\mathbb{R}^{n})$ 中也是 Cauchy 列. 由于 $W^{k-l,r}(\mathbb{R}^{n})$ 完备, 存在 $w\in W^{k-l,r}(\mathbb{R}^{n})$ 使得
\begin{align*}
u_{m}\to w\quad\text{于}\ W^{k-l,r}(\mathbb{R}^{n}).
\end{align*}
又由 \eqref{eq::jijwIIJEI9-2-m}式知 $u_{m}\to\overline{u}$ 于 $W^{k,p}(\mathbb{R}^{n})$, 特别地 $u_{m}\to\overline{u}$ 于 $L^{p}(\mathbb{R}^{n})$. 另一方面, $u_{m}\to w$ 于 $W^{k-l,r}(\mathbb{R}^{n})$, 特别地 $u_{m}\to w$ 于 $L^{r}(\mathbb{R}^{n})$. 因为 $r>n>p$ 且所有 $u_{m}$ 的支集都包含于某一固定紧集 $K$ 中, 由\cref{Real Analysis-proposition:Lp空间包含关系}, $L^{r}(K)$ 收敛蕴含 $L^{p}(K)$ 收敛, 即 $u_{m}\to w$ 于 $L^{p}(\mathbb{R}^{n})$. 由 $L^{p}$ 极限的唯一性, 得到 $w=\overline{u}$ a.e. 于 $\mathbb{R}^{n}$. 因此
\begin{align*}
\overline{u}\in W^{k-l,r}(\mathbb{R}^{n}),
\end{align*}
且在 \eqref{eq::jijwIIJEI9-2-l}式中令 $m\to\infty$, 得到
\begin{align}\label{eq::jijwIIJEI9-2-bar}
\|\overline{u}\|_{W^{k-l,r}(\mathbb{R}^{n})}\leqslant C_{l}'\|\overline{u}\|_{W^{k,p}(\mathbb{R}^{n})}.
\end{align}

对任意多重指标 $|\alpha|\leqslant k-l-1$, 由 $\overline{u}\in W^{k-l,r}(\mathbb{R}^{n})$ 知 $D^{\alpha}\overline{u}\in W^{1,r}(\mathbb{R}^{n})$. 因为 $r>n$, 由\hyperref[theorem:Morrey不等式]{Morrey不等式}, 对每个 $u_{m}$ 有
\begin{align}\label{eq::jijwIIJEI9-2-morrey}
\|D^{\alpha}u_{m}\|_{C^{0,1-\frac{n}{r}}(\mathbb{R}^{n})}\leqslant C\|D^{\alpha}u_{m}\|_{W^{1,r}(\mathbb{R}^{n})}\leqslant C\|u_{m}\|_{W^{k-l,r}(\mathbb{R}^{n})}\leqslant CC_{l}'\|u_{m}\|_{W^{k,p}(\mathbb{R}^{n})}.
\end{align}
由 \eqref{eq::jijwIIJEI9-2-cauchy}式与 \hyperref[theorem:Morrey不等式]{Morrey不等式}, 对任意 $m,m'\geqslant1$ 有
\begin{align*}
\|D^{\alpha}u_{m}-D^{\alpha}u_{m'}\|_{C^{0,1-\frac{n}{r}}(\mathbb{R}^{n})}\leqslant C\|D^{\alpha}u_{m}-D^{\alpha}u_{m'}\|_{W^{1,r}(\mathbb{R}^{n})}\leqslant C\|u_{m}-u_{m'}\|_{W^{k-l,r}(\mathbb{R}^{n})}.
\end{align*}
因此 $\{D^{\alpha}u_{m}\}$ 在 $C^{0,1-\frac{n}{r}}(\mathbb{R}^{n})$ 中为 Cauchy 列. 由于 $C^{0,1-\frac{n}{r}}(\mathbb{R}^{n})$ 完备, 存在 $v_{\alpha}\in C^{0,1-\frac{n}{r}}(\mathbb{R}^{n})$ 使得
\begin{align*}
D^{\alpha}u_{m}\to v_{\alpha}\quad\text{于}\ C^{0,1-\frac{n}{r}}(\mathbb{R}^{n}).
\end{align*}
另一方面, 由 $u_{m}\to\overline{u}$ 于 $W^{k-l,r}(\mathbb{R}^{n})$ 知 $D^{\alpha}u_{m}\to D^{\alpha}\overline{u}$ 于$L^{r}(\mathbb{R}^{n})$. 而 $C^{0,1-\frac{n}{r}}$ 收敛蕴含一致收敛, 从而蕴含 $L^{r}(\mathbb{R}^{n})$ 收敛. 由 $L^{r}$ 极限的唯一性, 得到 $v_{\alpha}=D^{\alpha}\overline{u}$ a.e. 于 $\mathbb{R}^{n}$. 因此
\begin{align*}
D^{\alpha}\overline{u}\in C^{0,1-\frac{n}{r}}(\mathbb{R}^{n}),
\end{align*}
且在 \eqref{eq::jijwIIJEI9-2-morrey}式中令 $m\to\infty$, 由范数的连续性得到
\begin{align}\label{eq::jijwIIJEI9-2-holder}
\|D^{\alpha}\overline{u}\|_{C^{0,1-\frac{n}{r}}(\mathbb{R}^{n})}\leqslant C\|\overline{u}\|_{W^{k-l,r}(\mathbb{R}^{n})}\leqslant CC_{l}'\|\overline{u}\|_{W^{k,p}(\mathbb{R}^{n})}.
\end{align}
注意到
\begin{align*}
1-\frac{n}{r}=1-\frac{n}{p}+l=\left[\frac{n}{p}\right]+1-\frac{n}{p}.
\end{align*}
因此
\begin{align*}
\overline{u}\in C^{k-\left[\frac{n}{p}\right]-1,\left[\frac{n}{p}\right]+1-\frac{n}{p}}(\mathbb{R}^{n}),
\end{align*}
并且由 \eqref{eq::jijwIIJEI9-2-holder}式得到
\begin{align}\label{eq::jijwIIJEI9-2-final}
\|\overline{u}\|_{C^{k-\left[\frac{n}{p}\right]-1,\left[\frac{n}{p}\right]+1-\frac{n}{p}}(\mathbb{R}^{n})}\leqslant C\|\overline{u}\|_{W^{k,p}(\mathbb{R}^{n})}.
\end{align}

最后利用扩张估计 \eqref{eq::jijwIIJEI9-2}, 便得到
\begin{align*}
\left\| u \right\| _{C^{k-\left[ \frac{n}{p} \right] -1,\left[ \frac{n}{p} \right] +1-\frac{n}{p}}(\overline{U})}\leqslant \left\| \overline{u} \right\| _{C^{k-\left[ \frac{n}{p} \right] -1,\left[ \frac{n}{p} \right] +1-\frac{n}{p}}(\mathbb{R} ^n)}\leqslant C\left\| \overline{u} \right\| _{W^{k,p}(\mathbb{R} ^n)}\leqslant CC_0\left\| u \right\| _{W^{k,p}(U)}.
\end{align*}
上述常数 $CC_0$ 由每一步的常数复合而成, 只依赖于 $k$、$p$、$n$ 与 $U$.

3. 最后设 $k>\frac{n}{p}$ 成立且 $\frac{n}{p}$ 为整数. 令
\begin{align*}
m\triangleq\frac{n}{p}\in\mathbb{N},\qquad l\triangleq m-1=\frac{n}{p}-1.
\end{align*}
由于 $\partial U$ 为 $C^{1}$, 由\hyperref[theorem:ch05-扩张定理]{扩张定理}, 存在 $\overline{u}\in W^{k,p}(\mathbb{R}^{n})$, 使得
\begin{align}\label{eq::jijwIIJEI9-3}
\overline{u}=u\ \text{于}\ U,\qquad \operatorname{supp}\overline{u}\ \text{紧},\qquad
\|\overline{u}\|_{W^{k,p}(\mathbb{R}^{n})}\leqslant C_{0}\|u\|_{W^{k,p}(U)}.
\end{align}
由\hyperref[theorem:ch05-光滑函数的局部逼近定理]{光滑函数的局部逼近定理}, 取开集为 $\mathbb{R}^{n}$, 存在 $u_{m}\in C_{c}^{\infty}(\mathbb{R}^{n})$, 使得
\begin{align}\label{eq::jijwIIJEI9-3-m}
u_{m}\to\overline{u}\quad\text{于}\ W^{k,p}(\mathbb{R}^{n}).
\end{align}
特别地, $\{u_{m}\}$ 在 $W^{k,p}(\mathbb{R}^{n})$ 中为 Cauchy 列.
先对 $u_{m}$ 证明相应估计. 令
\begin{align*}
p_{0}\triangleq p,\qquad \frac{1}{p_{j+1}}\triangleq\frac{1}{p_{j}}-\frac{1}{n},\quad j=0,1,\cdots,l-1.
\end{align*}
由于 $l=m-1<\frac{n}{p}=m$, 故对每个 $j=0,1,\cdots,l-1$, 有 $p_{j}<n$, 从而\hyperref[theorem:Gagliardo--Nirenberg--Sobolev不等式]{Gagliardo--Nirenberg--Sobolev不等式}可以应用.

对任意多重指标 $|\beta|\leqslant k-1$, 令 $v\triangleq D^{\beta}u_{m}$. 由于 $u_{m}\in C_{c}^{\infty}(\mathbb{R}^{n})$, 故 $v\in C_{c}^{\infty}(\mathbb{R}^{n})$. 对 $v$ 应用\hyperref[theorem:Gagliardo--Nirenberg--Sobolev不等式]{Gagliardo--Nirenberg--Sobolev不等式}, 得到
\begin{align}\label{eq::jijwIIJEI9-3-first}
\|D^{\beta}u_{m}\|_{L^{p_{1}}(\mathbb{R}^{n})}
=\|v\|_{L^{(p_{0})^{*}}(\mathbb{R}^{n})}\leqslant C_{1}\|Dv\|_{L^{p_{0}}(\mathbb{R}^{n})}=C_{1}\|D(D^{\beta}u_{m})\|_{L^{p}(\mathbb{R}^{n})}\leqslant C_{1}\|u_{m}\|_{W^{k,p}(\mathbb{R}^{n})}.
\end{align}
因此 $u_{m}\in W^{k-1,p_{1}}(\mathbb{R}^{n})$, 且
\begin{align*}
\left\| u_m \right\| _{W^{k-1,p_1}\left( \mathbb{R} ^n \right)}&=\left( \sum_{\left| \beta \right|\leqslant k-1}{\left\| D^{\beta}u_m \right\| _{L^{p_1}\left( \mathbb{R} ^n \right)}^{p_1}} \right) ^{\frac{1}{p_1}}\leqslant \left( \sum_{\left| \beta \right|\leqslant k-1}{C_1^{p_1}\left\| u_m \right\| _{W^{k,p}\left( \mathbb{R} ^n \right)}^{p_1}} \right) ^{\frac{1}{p_1}}
\\&=C_{1}\left( \sum_{\left| \beta \right|\leqslant k-1}{1} \right) ^{\frac{1}{p_1}}\left\| u_m \right\| _{W^{k,p}\left( \mathbb{R} ^n \right)}\triangleq C_{1}^{\prime}\left\| u_m \right\| _{W^{k,p}\left( \mathbb{R} ^n \right)}.
\end{align*}
假设对某个 $j\in\{0,1,\cdots,l-1\}$, 已有
\begin{align*}
u_{m}\in W^{k-j,p_{j}}(\mathbb{R}^{n}),
\end{align*}
且
\begin{align*}
\|u_{m}\|_{W^{k-j,p_{j}}(\mathbb{R}^{n})}\leqslant C_{j}'\|u_{m}\|_{W^{k,p}(\mathbb{R}^{n})}.
\end{align*}
对任意多重指标 $|\beta|\leqslant k-j-1$, 令 $v\triangleq D^{\beta}u_{m}$. 因为 $u_{m}\in W^{k-j,p_{j}}(\mathbb{R}^{n})$, 所以 $D^{\beta}u_{m}\in W^{1,p_{j}}(\mathbb{R}^{n})$, 且 $v\in C_{c}^{\infty}(\mathbb{R}^{n})$. 对 $v$ 应用\hyperref[theorem:Gagliardo--Nirenberg--Sobolev不等式]{Gagliardo--Nirenberg--Sobolev不等式}, 得到
\begin{align*}
\|D^{\beta}u_{m}\|_{L^{p_{j+1}}(\mathbb{R}^{n})}
&=\|v\|_{L^{(p_{j})^{*}}(\mathbb{R}^{n})}\leqslant C_{j+1}\|Dv\|_{L^{p_{j}}(\mathbb{R}^{n})}=C_{j+1}\|D(D^{\beta}u_{m})\|_{L^{p_{j}}(\mathbb{R}^{n})}\\
&\leqslant C_{j+1}\|u_{m}\|_{W^{k-j,p_{j}}(\mathbb{R}^{n})}\leqslant C_{j+1}C_{j}'\|u_{m}\|_{W^{k,p}(\mathbb{R}^{n})}.
\end{align*}
因此 $u_{m}\in W^{k-j-1,p_{j+1}}(\mathbb{R}^{n})$, 且
\begin{align*}
\left\| u_m \right\| _{W^{k-j-1,p_{j+1}}\left( \mathbb{R} ^n \right)}
&=\left( \sum_{\left| \beta \right|\leqslant k-j-1}{\left\| D^{\beta}u_m \right\| _{L^{p_{j+1}}\left( \mathbb{R} ^n \right)}^{p_{j+1}}} \right) ^{\frac{1}{p_{j+1}}}\leqslant \left( \sum_{\left| \beta \right|\leqslant k-j-1}{\left( C_{j+1}C_{j}^{\prime}\left\| u_m \right\| _{W^{k,p}\left( \mathbb{R} ^n \right)} \right) ^{p_{j+1}}} \right) ^{\frac{1}{p_{j+1}}}\\
&=C_{j+1}C_{j}^{\prime}\left( \sum_{\left| \beta \right|\leqslant k-j-1}{1} \right) ^{\frac{1}{p_{j+1}}}\left\| u_m \right\| _{W^{k,p}\left( \mathbb{R} ^n \right)}\triangleq C_{j+1}^{\prime}\left\| u_m \right\| _{W^{k,p}\left( \mathbb{R} ^n \right)}.
\end{align*}

重复上述论证, 经过 $l$ 步后，即取$j=l-1$,得到
\begin{align*}
u_{m}\in W^{k-l,p_{l}}(\mathbb{R}^{n}),
\end{align*}
其中
\begin{align*}
\frac{1}{p_{l}}=\frac{1}{p}-\frac{l}{n}=\frac{m}{n}-\frac{m-1}{n}=\frac{1}{n},
\end{align*}
即 $p_{l}=n$. 因此
\begin{align*}
u_{m}\in W^{k-l,n}(\mathbb{R}^{n})=W^{k-m+1,n}(\mathbb{R}^{n}),
\end{align*}
且
\begin{align}\label{eq::jijwIIJEI9-3-l}
\|u_{m}\|_{W^{k-m+1,n}(\mathbb{R}^{n})}\leqslant C\|u_{m}\|_{W^{k,p}(\mathbb{R}^{n})}.
\end{align}

由 \eqref{eq::jijwIIJEI9-3-m}式知 $u_{m}\to\overline{u}$ 于 $W^{k,p}(\mathbb{R}^{n})$, 故 $\{u_{m}\}$ 在 $W^{k,p}(\mathbb{R}^{n})$ 中为 Cauchy 列. 由\eqref{eq::jijwIIJEI9-3-l}式知, 对任意 $m,m'\geqslant1$ 有
\begin{align}\label{eq::jijwIIJEI9-3-cauchy}
\|u_{m}-u_{m'}\|_{W^{k-m+1,n}(\mathbb{R}^{n})}\leqslant C\|u_{m}-u_{m'}\|_{W^{k,p}(\mathbb{R}^{n})}.
\end{align}
因此 $\{u_{m}\}$ 在 $W^{k-m+1,n}(\mathbb{R}^{n})$ 中也是 Cauchy 列. 由于 $W^{k-m+1,n}(\mathbb{R}^{n})$ 完备, 存在 $w\in W^{k-m+1,n}(\mathbb{R}^{n})$ 使得
\begin{align*}
u_{m}\to w\quad\text{于}\ W^{k-m+1,n}(\mathbb{R}^{n}).
\end{align*}
又由 \eqref{eq::jijwIIJEI9-3-m}式知 $u_{m}\to\overline{u}$ 于 $W^{k,p}(\mathbb{R}^{n})$, 特别地 $u_{m}\to\overline{u}$ 于 $L^{p}(\mathbb{R}^{n})$. 另一方面, $u_{m}\to w$ 于 $W^{k-m+1,n}(\mathbb{R}^{n})$, 特别地 $u_{m}\to w$ 于 $L^{n}(\mathbb{R}^{n})$. 因为$\frac{n}{p}>0$且是整数,所以$n\geqslant p$且所有 $u_{m}$ 的支集都包含于某一固定紧集 $K$ 中, 由\cref{Real Analysis-proposition:Lp空间包含关系}, $L^{n}(K)$ 收敛蕴含 $L^{p}(K)$ 收敛, 即 $u_{m}\to w$ 于 $L^{p}(\mathbb{R}^{n})$. 由 $L^{p}$ 极限的唯一性, 得到 $w=\overline{u}$ a.e. 于 $\mathbb{R}^{n}$. 因此
\begin{align*}
\overline{u}\in W^{k-m+1,n}(\mathbb{R}^{n}),
\end{align*}
且在 \eqref{eq::jijwIIJEI9-3-l}式中令 $m\to\infty$, 得到
\begin{align}\label{eq::jijwIIJEI9-3-bar}
\|\overline{u}\|_{W^{k-m+1,n}(\mathbb{R}^{n})}\leqslant C\|\overline{u}\|_{W^{k,p}(\mathbb{R}^{n})}.
\end{align}

对任意多重指标 $|\alpha|\leqslant k-m$, 由 $\overline{u}\in W^{k-m+1,n}(\mathbb{R}^{n})$ 知 $D^{\alpha}\overline{u}\in W^{1,n}(\mathbb{R}^{n})$. 由\cref{corollary:Sobolev嵌入Lqq大于}, 对每个 $u_{m}$ 及任意 $n\leqslant q<\infty$ 有
\begin{align}\label{eq::jijwIIJEI9-3-sob}
\|D^{\alpha}u_{m}\|_{L^{q}(\mathbb{R}^{n})}\leqslant C\|D^{\alpha}u_{m}\|_{W^{1,n}(\mathbb{R}^{n})}\leqslant C\|u_{m}\|_{W^{k-m+1,n}(\mathbb{R}^{n})}\leqslant C\|u_{m}\|_{W^{k,p}(\mathbb{R}^{n})}.
\end{align}
由 \eqref{eq::jijwIIJEI9-3-cauchy}式与上式, 对任意 $m,m'\geqslant1$ 有
\begin{align*}
\|D^{\alpha}u_{m}-D^{\alpha}u_{m'}\|_{L^{q}(\mathbb{R}^{n})}\leqslant C\|u_{m}-u_{m'}\|_{W^{k-m+1,n}(\mathbb{R}^{n})}.
\end{align*}
因此 $\{D^{\alpha}u_{m}\}$ 在 $L^{q}(\mathbb{R}^{n})$ 中为 Cauchy 列. 由于 $L^{q}(\mathbb{R}^{n})$ 完备, 存在 $v_{\alpha}\in L^{q}(\mathbb{R}^{n})$ 使得
\begin{align*}
D^{\alpha}u_{m}\to v_{\alpha}\quad\text{于}\ L^{q}(\mathbb{R}^{n}).
\end{align*}
另一方面, 由 $u_{m}\to\overline{u}$ 于 $W^{k-m+1,n}(\mathbb{R}^{n})$ 知 $D^{\alpha}u_{m}\to D^{\alpha}\overline{u}$ 于 $L^{n}(\mathbb{R}^{n})$. 因 $q\geqslant n$ 且支集有界,故由\cref{Real Analysis-proposition:Lp空间包含关系}知$L^{q}$ 收敛蕴含 $L^{n}$ 收敛.由 $L^{n}$ 极限的唯一性得到 $v_{\alpha}=D^{\alpha}\overline{u}$ a.e. 于 $\mathbb{R}^{n}$. 因此
\begin{align*}
D^{\alpha}\overline{u}\in L^{q}(\mathbb{R}^{n}),
\end{align*}
且在 \eqref{eq::jijwIIJEI9-3-sob}式中令 $m\to\infty$, 得到
\begin{align}\label{eq::jijwIIJEI9-3-Lq}
\|D^{\alpha}\overline{u}\|_{L^{q}(\mathbb{R}^{n})}\leqslant C\|\overline{u}\|_{W^{k,p}(\mathbb{R}^{n})},\quad n\leqslant q<\infty.
\end{align}

现在对任意 $|\alpha|\leqslant k-m-1$, 由 \eqref{eq::jijwIIJEI9-3-Lq}式知 $D^{\alpha}\overline{u}\in W^{1,q}(\mathbb{R}^{n})$ 对任意 $n\leqslant q<\infty$ 成立. 取 $q>n$, 由\hyperref[theorem:Morrey不等式]{Morrey不等式}, 对每个 $u_{m}$ 有
\begin{align}\label{eq::jijwIIJEI9-3-morrey}
\|D^{\alpha}u_{m}\|_{C^{0,1-\frac{n}{q}}(\mathbb{R}^{n})}\leqslant C\|D^{\alpha}u_{m}\|_{W^{1,q}(\mathbb{R}^{n})}\leqslant C\|u_{m}\|_{W^{k,p}(\mathbb{R}^{n})}.
\end{align}
由\eqref{eq::jijwIIJEI9-3-cauchy}式与 \hyperref[theorem:Morrey不等式]{Morrey不等式}, 对任意 $m,m'\geqslant1$ 有
\begin{align*}
\|D^{\alpha}u_{m}-D^{\alpha}u_{m'}\|_{C^{0,1-\frac{n}{q}}(\mathbb{R}^{n})}\leqslant C\|u_{m}-u_{m'}\|_{W^{k-m+1,n}(\mathbb{R}^{n})}.
\end{align*}
因此 $\{D^{\alpha}u_{m}\}$ 在 $C^{0,1-\frac{n}{q}}(\mathbb{R}^{n})$ 中为 Cauchy 列. 由于 $C^{0,1-\frac{n}{q}}(\mathbb{R}^{n})$ 完备, 存在 $v_{\alpha}\in C^{0,1-\frac{n}{q}}(\mathbb{R}^{n})$ 使得
\begin{align*}
D^{\alpha}u_{m}\to v_{\alpha}\quad\text{于}\ C^{0,1-\frac{n}{q}}(\mathbb{R}^{n}).
\end{align*}
另一方面, $D^{\alpha}u_{m}\to D^{\alpha}\overline{u}$ 于 $L^{q}(\mathbb{R}^{n})$, 且 $C^{0,1-\frac{n}{q}}$ 收敛蕴含一致收敛, 从而蕴含 $L^{q}(\mathbb{R}^{n})$ 收敛. 由 $L^{q}$ 极限的唯一性, 得到 $v_{\alpha}=D^{\alpha}\overline{u}$ a.e. 于 $\mathbb{R}^{n}$. 因此
\begin{align*}
D^{\alpha}\overline{u}\in C^{0,1-\frac{n}{q}}(\mathbb{R}^{n}),
\end{align*}
且在 \eqref{eq::jijwIIJEI9-3-morrey}式中令 $m\to\infty$, 得到
\begin{align}\label{eq::jijwIIJEI9-3-holder}
\|D^{\alpha}\overline{u}\|_{C^{0,1-\frac{n}{q}}(\mathbb{R}^{n})}\leqslant C\|\overline{u}\|_{W^{k,p}(\mathbb{R}^{n})},\quad q>n.
\end{align}
因为 $q>n$ 可以任意大, 所以对任意 $0<\gamma<1$, 取 $q$ 充分大使得 $1-\frac{n}{q}>\gamma$, 就有
\begin{align*}
D^{\alpha}\overline{u}\in C^{0,1-\frac{n}{q}}(\mathbb{R}^{n})\subset C^{0,\gamma}(\mathbb{R}^{n}).
\end{align*}
因此
\begin{align*}
\overline{u}\in C^{k-m-1,\gamma}(\mathbb{R}^{n}),
\end{align*}
并且由 \eqref{eq::jijwIIJEI9-3-holder}式得到
\begin{align}\label{eq::jijwIIJEI9-3-final}
\|\overline{u}\|_{C^{k-m-1,\gamma}(\mathbb{R}^{n})}\leqslant C\|\overline{u}\|_{W^{k,p}(\mathbb{R}^{n})}.
\end{align}

最后利用扩张估计 \eqref{eq::jijwIIJEI9-3}式, 便得到
\begin{align*}
\left\| u \right\| _{C^{k-m-1,\gamma}(\overline{U})}\leqslant \left\| \overline{u} \right\| _{C^{k-m-1,\gamma}(\mathbb{R} ^n)}\leqslant C\left\| \overline{u} \right\| _{W^{k,p}(\mathbb{R} ^n)}\leqslant CC_0\left\| u \right\| _{W^{k,p}(U)}.
\end{align*}
上述常数 $CC_0$ 由每一步的常数复合而成, 只依赖于 $k$、$p$、$n$、$\gamma$ 与 $U$.
\end{proof}

各种一般的Sobolev型不等式也可利用Fourier变换证明：见习题 20。

\end{document}
